\chapter{Complexity — Nondeterminism}
%%% generated by the blueprint pipeline from TCSlib.Complexity.ClassNP.Nondeterminism
%%%%%%%%%%%%%%%%%%%%%%%%%%%%%%%%%%%%%%%%%%%%%%%%%%%%%%%%%%%%%%%%%%%%%%%%%%%%%%%
\section{Overview}

This module proves the nondeterministic-machine characterizations of the
certificate classes from Arora--Barak: Theorem~2.6, stating that
$\mathrm{NP} = \bigcup_c \mathrm{NTIME}(n^c + 1)$ (with the $+1$ padding recorded at
the definition of $\mathrm{P}$); the analogous identity
$\mathrm{NEXP} = \bigcup_c \mathrm{NTIME}(2^{n^c})$, which reconciles the
verifier-certificate form of $\mathrm{NEXP}$ with its textbook machine definition;
and Theorem~2.22, that $\mathrm{EXP} \ne \mathrm{NEXP}$ implies
$\mathrm{P} \ne \mathrm{NP}$, by padding. Both directions of each characterization
are carried out as explicit machine compilations: a captured choice-word simulator
turns a nondeterministic decider into a polynomial-time verifier with exact-length
certificates, while a scheduler-driven reverse host guesses a certificate of the
exact scheduled length and runs the relocated verifier, with every phase proved as
a timed configuration contract.

%%%%%%%%%%%%%%%%%%%%%%%%%%%%%%%%%%%%%%%%%%%%%%%%%%%%%%%%%%%%%%%%%%%%%%%%%%%%%%%
\section{Declarations}

\subsection{Polynomial split and the choice-word verifier}

\begin{lemma}[Strict monotonicity of the padded length map]
\lean{Complexity.certificate_split_strictMono}
\label{Complexity.certificate_split_strictMono}
\leanok
\difficulty{2}
For all natural numbers $C$ and $c$, the map $n \mapsto n + C\,(n+1)^c$ on the natural numbers is strictly increasing.
\end{lemma}

\begin{lemma}[Uniqueness of the certificate split]
\lean{Complexity.certificate_split_unique}
\label{Complexity.certificate_split_unique}
\leanok
\uses{Complexity.certificateSplit, Complexity.certificate_split_strictMono}
\difficulty{3}
Let $C, c$ be natural numbers and let $x, u, y, v$ be binary strings with $|u| = C\,(|x|+1)^c$ and $|v| = C\,(|y|+1)^c$. If the concatenations $x\,u$ and $y\,v$ are equal as strings, then $x = y$ and $u = v$.
\end{lemma}

\begin{lemma}[Failed split search characterizes unsolvable lengths]
\lean{Complexity.certificateSplit_none_iff}
\label{Complexity.certificateSplit_none_iff}
\leanok
\uses{Complexity.certificateSplit}
\difficulty{3}
For natural numbers $C$, $c$ and $m$, the bounded search for a split position (which scans all $n \le m$ for a solution of $n + C\,(n+1)^c = m$) returns no result if and only if no natural number $n$ satisfies $n + C\,(n+1)^c = m$.
\end{lemma}

\begin{lemma}[Completeness of the split search]
\lean{Complexity.certificateSplit_complete}
\label{Complexity.certificateSplit_complete}
\leanok
\uses{Complexity.certificateSplit, Complexity.certificateSplit_none_iff, Complexity.certificateSplit_spec, Complexity.certificate_split_strictMono}
\difficulty{3}
For natural numbers $C$, $c$, $m$ and $n$ with $n + C\,(n+1)^c = m$, the bounded search for a split position at length $m$ succeeds and returns exactly $n$.
\end{lemma}

\begin{definition}[The choice-word verifier language]
\lean{Complexity.choiceVerifier}
\label{Complexity.choiceVerifier}
\leanok
\uses{Turing.FinNDTM, Turing.NDTM.initCfg, Turing.NDTM.runWith}
Given a nondeterministic Turing machine $N$ over the binary alphabet and natural numbers $C$ and $c$, the choice-word verifier language of $N$ consists of all binary strings of the form $x\,u$ (concatenation) where $|u| = C\,(|x|+1)^c$ and the run of $N$ on input $x$ under choice word $u$ has halted with output exactly the one-bit string $1$.
\end{definition}

\begin{lemma}[Membership of a correctly split word]
\lean{Complexity.choiceVerifier_append}
\label{Complexity.choiceVerifier_append}
\leanok
\uses{Complexity.certificate_split_unique, Complexity.choiceVerifier, Turing.FinNDTM, Turing.NDTM.initCfg, Turing.NDTM.runWith}
\difficulty{3}
Let $N$ be a nondeterministic Turing machine over the binary alphabet, let $C, c$ be natural numbers, and let $x$ and $u$ be binary strings with $|u| = C\,(|x|+1)^c$. Then the concatenation $x\,u$ belongs to the choice-word verifier language of $N$ with parameters $C, c$ if and only if the run of $N$ on input $x$ under choice word $u$ has halted with output exactly the one-bit string $1$.
\end{lemma}

\begin{lemma}[Rejection of words without a valid split]
\lean{Complexity.choiceVerifier_no_split}
\label{Complexity.choiceVerifier_no_split}
\leanok
\uses{Complexity.certificateSplit, Complexity.certificateSplit_none_iff, Complexity.choiceVerifier, Turing.FinNDTM}
\difficulty{3}
Let $N$ be a nondeterministic Turing machine over the binary alphabet and let $C, c$ be natural numbers. If no natural number $n$ satisfies $n + C\,(n+1)^c = |y|$ for a binary string $y$ (so the bounded split search fails at length $|y|$), then $y$ does not belong to the choice-word verifier language of $N$ with parameters $C, c$.
\end{lemma}

\begin{lemma}[Budget invariance of acceptance after all branches halt]
\lean{Complexity.acceptsWithin_iff_of_halts}
\label{Complexity.acceptsWithin_iff_of_halts}
\leanok
\uses{Turing.FinNDTM, Turing.FinNDTM.AcceptsWithin, Turing.FinNDTM.AcceptsWithin.mono, Turing.NDTM.HaltsWithin, Turing.NDTM.initCfg, Turing.NDTM.runWith_append, Turing.NDTM.runWith_of_halt}
\difficulty{4}
Let $N$ be a nondeterministic Turing machine over the binary alphabet, let $x$ be an input string, and let $t \le t'$ be step budgets such that every choice word of length $t$ drives $N$ on $x$ into a halted configuration. Then $N$ accepts $x$ within $t'$ steps if and only if $N$ accepts $x$ within $t$ steps, where acceptance within a budget means that some choice word of exactly that length yields a halted run with output the one-bit string $1$.
\end{lemma}

\begin{lemma}[The doubled coefficient covers the time bound]
\lean{Complexity.choice_budget_le}
\label{Complexity.choice_budget_le}
\leanok
\difficulty{3}
For all natural numbers $a$, $c$ and $n$, the inequality $a\,(n^c + 1) \le 2a\,(n+1)^c$ holds.
\end{lemma}

\begin{lemma}[Certificates for the choice-word verifier]
\lean{Complexity.choice_certificate_iff}
\label{Complexity.choice_certificate_iff}
\leanok
\uses{Complexity.acceptsWithin_iff_of_halts, Complexity.capturedSummary, Complexity.choiceVerifier, Complexity.choiceVerifier_append, Complexity.choice_budget_le, Turing.FinNDTM, Turing.FinNDTM.DecidesInTime}
\difficulty{3}
Let $N$ be a nondeterministic Turing machine over the binary alphabet that decides a language $L$ within time $a\,(n^c+1)$ on inputs of length $n$, meaning that every branch halts within the budget and acceptance characterizes membership. Then for every binary string $x$: $x \in L$ if and only if there is a binary string $u$ with $|u| = 2a\,(|x|+1)^c$ such that the concatenation $x\,u$ lies in the choice-word verifier language of $N$ with parameters $2a$ and $c$.
\end{lemma}

\subsection{Output capture and the deterministic simulation core}

\begin{definition}[Finite summary of a captured output word]
\lean{Complexity.capturedSummary}
\label{Complexity.capturedSummary}
\leanok
For a binary string $w$, the output summary is: nothing when $w$ is empty, the bit $b$ when $w$ is the single-bit string $b$, and the bit $0$ when $w$ has two or more bits. It is a three-valued digest that recognizes exactly the singleton outputs.
\end{definition}

\begin{definition}[One-step update of the output summary]
\lean{Complexity.captureEmission}
\label{Complexity.captureEmission}
\leanok
Given a current output summary $s$ (an optional bit) and an optional emitted bit $e$, the updated summary is $s$ when nothing is emitted, the emitted bit when $s$ is still empty, and the bit $0$ when a bit is emitted while $s$ already holds one.
\end{definition}

\begin{lemma}[The summary tracks appended emissions]
\lean{Complexity.captureEmission_correct}
\label{Complexity.captureEmission_correct}
\leanok
\uses{Complexity.captureEmission, Complexity.capturedSummary}
\difficulty{2}
For every binary string $w$ and every optional bit $e$, updating the output summary of $w$ by the emission $e$ equals the output summary of $w$ with the content of $e$ appended.
\end{lemma}

\begin{lemma}[The accepting summary characterizes the singleton]
\lean{Complexity.capturedSummary_true}
\label{Complexity.capturedSummary_true}
\leanok
\uses{Complexity.capturedSummary}
\difficulty{2}
For every binary string $w$, the output summary of $w$ equals the bit $1$ if and only if $w$ is exactly the one-bit string $1$.
\end{lemma}

\begin{definition}[Partition of work tapes into source and private blocks]
\lean{Complexity.choiceTapes}
\label{Complexity.choiceTapes}
\leanok
Given an assignment of values to $k$ source tapes and three further values, the combined assignment on $k+3$ tapes is the function that places the source values on the first $k$ tapes and the three extra values (a virtual input, a choice buffer, and a captured output) on the last three tapes, in that order.
\end{definition}

\begin{definition}[Deterministic simulation core for a fixed NDTM]
\lean{Complexity.choiceCore}
\label{Complexity.choiceCore}
\leanok
\uses{Complexity.captureEmission, Complexity.choiceTapes, Turing.FinNDTM, Turing.FinTM, Turing.FinTM.virtualMove, Turing.FinTM.virtualNextTag}
For a nondeterministic Turing machine $N$ over the binary alphabet with $k$ work tapes, a deterministic machine on $k+3$ tapes is defined that simulates $N$ from prepared private tapes: each bit read from the choice buffer drives one step of $N$ (with the input read from a copied virtual-input tape and every emission written to a capture tape), and the first blank on the choice buffer triggers a final transition that halts and emits one verdict bit, which is $1$ exactly when the simulated machine has halted with accepting output summary.
\end{definition}

\begin{definition}[Embedding a source configuration in the simulation core]
\lean{Complexity.choiceCoreCfg}
\label{Complexity.choiceCoreCfg}
\leanok
\uses{Complexity.capturedSummary, Complexity.choiceCore, Complexity.choiceTapes, Turing.Cfg, Turing.FinNDTM, Turing.FinTM.bufferTape}
Given a configuration of a nondeterministic machine $N$ on input $x$, a boundary tag, a choice word $u$, a choice-head index $j$, and a physical input position, the corresponding embedding is the configuration of the simulation core of $N$ that stores the source state, tag and output summary in finite control, each source work tape on its own tape, and the words $x$, $u$ and the source output on the three private buffer tapes, with empty physical output.
\end{definition}

\begin{lemma}[One simulated step consumes one copied choice]
\lean{Complexity.choiceCore_step}
\label{Complexity.choiceCore_step}
\leanok
\uses{Complexity.captureEmission, Complexity.captureEmission_correct, Complexity.capturedSummary, Complexity.choiceCore, Complexity.choiceCoreCfg, Complexity.choiceTapes, Turing.Action.apply, Turing.Cfg, Turing.Cfg.inputSymbol, Turing.Cfg.workTapeSymbols, Turing.FinNDTM, Turing.FinTM.VirtualTag, Turing.FinTM.bufferTape_append, Turing.FinTM.bufferTape_inputSymbol, Turing.FinTM.virtualMove, Turing.FinTM.virtualMove_correct, Turing.FinTM.virtualNextTag, Turing.MultiTapeTM.step, Turing.NDTM.stepWith, Turing.NDTM.stepWith_of_halt, Turing.moveInputPos_zero}
\difficulty{6}
Let $N$ be a nondeterministic Turing machine over the binary alphabet, let a configuration of $N$ on input $x$ carry a valid virtual boundary tag, and let $u$ be a choice word with next unread index $j < |u|$. Then one physical step of the simulation core of $N$ transforms the embedded configuration into the embedding of the configuration obtained by one step of $N$ under the choice bit at position $j$, with choice index $j+1$ and some valid updated boundary tag.
\end{lemma}

\begin{lemma}[Lockstep simulation over a prefix of choices]
\lean{Complexity.choiceCore_run}
\label{Complexity.choiceCore_run}
\leanok
\uses{Complexity.choiceCore, Complexity.choiceCoreCfg, Complexity.choiceCore_step, Turing.Cfg, Turing.FinNDTM, Turing.FinTM.VirtualTag, Turing.MultiTapeTM.runFrom, Turing.MultiTapeTM.runFrom_succ_eq_step', Turing.NDTM.runWith, Turing.NDTM.runWith_append, Turing.NDTM.stepWith}
\difficulty{4}
Let $N$ be a nondeterministic Turing machine over the binary alphabet, let a configuration of $N$ on input $x$ carry a valid virtual boundary tag, and let $u$ be a choice word with $t \le |u|$. Then $t$ physical steps of the simulation core of $N$ carry the embedded configuration to the embedding of the run of $N$ under the first $t$ bits of $u$, again with a valid boundary tag.
\end{lemma}

\begin{lemma}[Final verdict transition at the choice blank]
\lean{Complexity.choiceCore_finish}
\label{Complexity.choiceCore_finish}
\leanok
\uses{Complexity.capturedSummary, Complexity.capturedSummary_true, Complexity.choiceCore, Complexity.choiceCoreCfg, Complexity.choiceTapes, Turing.Cfg, Turing.Cfg.workTapeSymbols, Turing.FinNDTM, Turing.MultiTapeTM.step}
\difficulty{3}
Let $N$ be a nondeterministic Turing machine over the binary alphabet and consider the embedding of a configuration of $N$ whose choice head stands at the blank after the complete choice word $u$. One further step of the simulation core halts with output a single bit, and that bit is $1$ exactly when the embedded source configuration has halted with output the one-bit string $1$.
\end{lemma}

\begin{lemma}[Timed contract of the simulation core]
\lean{Complexity.choiceCore_timed}
\label{Complexity.choiceCore_timed}
\leanok
\uses{Complexity.choiceCore, Complexity.choiceCoreCfg, Complexity.choiceCore_finish, Complexity.choiceCore_run, Turing.Cfg, Turing.FinNDTM, Turing.FinTM.VirtualTag, Turing.MultiTapeTM.runFrom, Turing.MultiTapeTM.runFrom_succ_eq_step', Turing.NDTM.runWith}
\difficulty{3}
Let $N$ be a nondeterministic Turing machine over the binary alphabet, let a configuration of $N$ carry a valid virtual boundary tag, and let $u$ be a choice word. From the prepared embedded configuration, the simulation core of $N$ halts after exactly $|u| + 1$ steps, and its output is the single bit recording whether the run of $N$ under $u$ from that configuration has halted with output the one-bit string $1$.
\end{lemma}

\begin{lemma}[Validity of the initial boundary tag]
\lean{Complexity.choiceCore_initial_tag}
\label{Complexity.choiceCore_initial_tag}
\leanok
\uses{Turing.Cfg.init, Turing.FinNDTM, Turing.FinTM.VirtualTag, Turing.NDTM.initCfg}
\difficulty{1}
For every nondeterministic Turing machine $N$ over the binary alphabet and every input string $x$, the initial configuration of $N$ on $x$ satisfies the virtual boundary-tag invariant with the interior tag value (its input head starts at position one, inside the window).
\end{lemma}

\subsection{The pair loader and the forward compilation}

\begin{lemma}[Agreement of the library and local split searches]
\lean{Complexity.cont_split_bridge}
\label{Complexity.cont_split_bridge}
\leanok
\uses{Complexity.certificateSplit, Complexity.contPairTM, Turing.solveSplit}
\difficulty{1}
For all natural numbers $C$, $c$ and $m$, the library split-search function and the local bounded split search return the same optional split position for the length equation $n + C\,(n+1)^c = m$.
\end{lemma}

\begin{definition}[Pair loader feeding the simulation core]
\lean{Complexity.contPairTM}
\label{Complexity.contPairTM}
\leanok
\uses{Complexity.choiceCore, Complexity.choiceTapes, Turing.Action.mapState, Turing.FinNDTM, Turing.FinTM}
For a nondeterministic Turing machine $N$ over the binary alphabet, a deterministic machine is defined that parses a self-delimiting encoded pair on its input tape: it decodes the doubled-bit prefix into an input buffer, copies the suffix after the separator into a choice buffer, rewinds both buffers, and then runs the simulation core of $N$ with its transition table embedded unchanged. An input that is not an aligned pair (in particular the empty word) is rejected with output bit $0$.
\end{definition}

\begin{definition}[Loader configuration of the pair machine]
\lean{Complexity.contLoadCfg}
\label{Complexity.contLoadCfg}
\leanok
\uses{Complexity.choiceTapes, Complexity.contPairTM, Turing.Cfg, Turing.FinNDTM, Turing.FinTM.bufferTape}
For a nondeterministic machine $N$, a loader state index, two buffer words $x$ and $u$, a physical input position and two buffer head positions, the corresponding configuration of the pair-loader machine of $N$ holds $x$ and $u$ on the two buffer tapes, with all source work tapes and the capture tape blank and the physical output empty.
\end{definition}

\begin{lemma}[Lockstep of the embedded core states]
\lean{Complexity.cont_core_run}
\label{Complexity.cont_core_run}
\leanok
\uses{Complexity.choiceCore, Complexity.contPairTM, Turing.Action.apply, Turing.Action.mapState, Turing.Cfg, Turing.Cfg.mapState, Turing.FinNDTM, Turing.MultiTapeTM.runFrom, Turing.MultiTapeTM.runFrom_comm_of_step, Turing.MultiTapeTM.step}
\difficulty{4}
Let $N$ be a nondeterministic machine over the binary alphabet and let $d$ be any configuration of the simulation core of $N$. For every step count $t$, running the pair-loader machine of $N$ for $t$ steps from the state-renamed copy of $d$ yields exactly the state-renamed copy of the $t$-step run of the simulation core from $d$.
\end{lemma}

\begin{lemma}[Exact cost of rewinding the choice buffer]
\lean{Complexity.cont_rewind_choices}
\label{Complexity.cont_rewind_choices}
\leanok
\uses{Complexity.choiceCoreCfg, Complexity.choiceTapes, Complexity.contLoadCfg, Complexity.contPairTM, Turing.Action.apply, Turing.Cfg.init, Turing.Cfg.mapState, Turing.Cfg.workTapeSymbols, Turing.FinNDTM, Turing.MultiTapeTM.runFrom, Turing.MultiTapeTM.runFrom_succ_eq_step, Turing.MultiTapeTM.runFrom_zero, Turing.MultiTapeTM.step, Turing.NDTM.initCfg, Turing.moveInputPos_zero}
\difficulty{5}
Let $N$ be a nondeterministic machine, let $x$ and $u$ be buffer words, and let $j \le |u|$. From the loader configuration in the choice-rewind state with choice head at position $j - 1$, exactly $j + 1$ steps of the pair-loader machine reach the renamed embedding of the initial configuration of $N$ on $x$ prepared with choice word $u$, preserving the physical input position.
\end{lemma}

\begin{lemma}[Exact cost of rewinding the input buffer]
\lean{Complexity.cont_rewind_input}
\label{Complexity.cont_rewind_input}
\leanok
\uses{Complexity.choiceTapes, Complexity.contLoadCfg, Complexity.contPairTM, Turing.Action.apply, Turing.Cfg.workTapeSymbols, Turing.FinNDTM, Turing.MultiTapeTM.runFrom, Turing.MultiTapeTM.runFrom_succ_eq_step, Turing.MultiTapeTM.runFrom_zero, Turing.MultiTapeTM.step, Turing.moveInputPos_zero}
\difficulty{5}
Let $N$ be a nondeterministic machine, let $x$ and $u$ be buffer words, and let $j \le |x|$. From the loader configuration in the input-rewind state with input-buffer head at position $j-1$ and choice head at the right end of $u$, exactly $j+1$ steps of the pair-loader machine reach the choice-rewind state with input-buffer head reset to zero and choice head at position $|u| - 1$, with all tape contents preserved.
\end{lemma}

\begin{lemma}[Buffer appends realized by one write action]
\lean{Complexity.cont_write_apply}
\label{Complexity.cont_write_apply}
\leanok
\uses{Complexity.choiceTapes, Complexity.contLoadCfg, Turing.Action.apply, Turing.FinNDTM, Turing.FinTM.bufferTape_append, Turing.moveInputPos}
\difficulty{4}
Let $N$ be a nondeterministic machine and consider a loader configuration whose two buffer heads stand at the right blanks of the buffer words $x$ and $u$. Applying an action that optionally writes one bit at each of the two buffer heads (advancing a head exactly when it writes) produces the loader configuration whose buffers are $x$ and $u$ with the respective optional bits appended, with the physical input head moved by the action's input move and all other tapes unchanged.
\end{lemma}

\begin{lemma}[Physical input read at an explicit position]
\lean{Complexity.cont_load_read}
\label{Complexity.cont_load_read}
\leanok
\uses{Complexity.contLoadCfg, Turing.Cfg.inputSymbol, Turing.FinNDTM, Turing.FinTM.inputSymbol_at}
\difficulty{1}
Let $N$ be a nondeterministic machine and consider a loader configuration over a physical input $y$ whose input head is at position $i + 1$ with $i \le |y|$. The symbol read from the physical input is the optional bit of $y$ at index $i$, which is blank exactly at the right boundary.
\end{lemma}

\begin{lemma}[One right-moving loader transition]
\lean{Complexity.cont_load_right}
\label{Complexity.cont_load_right}
\leanok
\uses{Complexity.choiceTapes, Complexity.contLoadCfg, Complexity.contPairTM, Complexity.cont_load_read, Complexity.cont_write_apply, Turing.Action.apply, Turing.Cfg.workTapeSymbols, Turing.FinNDTM, Turing.MultiTapeTM.step, Turing.moveInputPos_pos_of_ne_right}
\difficulty{4}
Let $N$ be a nondeterministic machine, let $y$ be a physical input with position index $i < |y|$, and suppose the transition table of the pair-loader machine, applied in a loader state $q$ to the bit of $y$ at index $i$, yields a right input move that optionally appends one bit to each buffer and enters a loader state $q'$. Then one physical step from the loader configuration at position $i+1$ with buffers $x$ and $u$ (heads at their right blanks) produces the loader configuration at position $i+2$ in state $q'$ with the appended buffers.
\end{lemma}

\begin{lemma}[Exact cost of copying the suffix]
\lean{Complexity.cont_copy_run}
\label{Complexity.cont_copy_run}
\leanok
\uses{Complexity.contLoadCfg, Complexity.contPairTM, Complexity.cont_load_right, Turing.FinNDTM, Turing.MultiTapeTM.runFrom, Turing.MultiTapeTM.runFrom_succ_eq_step, Turing.MultiTapeTM.runFrom_zero}
\difficulty{5}
Let $N$ be a nondeterministic machine and let the physical input $y$ decompose as $y = p\,r$ for words $p$ and $r$. From the suffix-copy loader state at position $|p| + 1$ with buffers $x$ and $u$, exactly $|r|$ steps of the pair-loader machine append $r$ to the choice buffer, ending at the right boundary of $y$ with the input buffer $x$ unchanged and no physical output.
\end{lemma}

\begin{lemma}[Entering the simulation core after copying]
\lean{Complexity.cont_start_core}
\label{Complexity.cont_start_core}
\leanok
\uses{Complexity.choiceCoreCfg, Complexity.choiceTapes, Complexity.contLoadCfg, Complexity.contPairTM, Complexity.cont_load_read, Complexity.cont_rewind_choices, Complexity.cont_rewind_input, Turing.Action.apply, Turing.Cfg.mapState, Turing.FinNDTM, Turing.MultiTapeTM.runFrom, Turing.MultiTapeTM.runFrom_add, Turing.MultiTapeTM.runFrom_succ_eq_step, Turing.MultiTapeTM.runFrom_zero, Turing.MultiTapeTM.step, Turing.NDTM.initCfg, Turing.moveInputPos_zero}
\difficulty{5}
Let $N$ be a nondeterministic machine and let $x$ and $u$ be the buffer words held by a loader configuration at the right boundary of the physical input $y$. Exactly $1 + (|x|+1) + (|u|+1)$ further steps of the pair-loader machine (one dispatch and the two buffer rewinds) reach the renamed embedding of the initial configuration of $N$ on $x$ prepared with choice word $u$; no bits of the choice word are consumed by these administrative steps.
\end{lemma}

\begin{lemma}[Composition of copying, rewinds and simulation]
\lean{Complexity.cont_suffix_run}
\label{Complexity.cont_suffix_run}
\leanok
\uses{Complexity.choiceCore_initial_tag, Complexity.choiceCore_timed, Complexity.contLoadCfg, Complexity.contPairTM, Complexity.cont_copy_run, Complexity.cont_core_run, Complexity.cont_start_core, Turing.Cfg.mapState, Turing.FinNDTM, Turing.MultiTapeTM.runFrom, Turing.MultiTapeTM.runFrom_add, Turing.NDTM.initCfg, Turing.NDTM.runWith}
\difficulty{4}
Let $N$ be a nondeterministic machine and let the physical input decompose as $y = p\,u$. Starting in the suffix-copy state at position $|p|+1$ with input buffer $x$ and empty choice buffer, after $|u| + (1 + (|x|+1) + (|u|+1)) + (|u|+1)$ steps the pair-loader machine has halted, and its output is the single bit recording whether the run of $N$ on $x$ under choice word $u$ has halted with output the one-bit string $1$.
\end{lemma}

\begin{lemma}[Two-step decoding of a doubled bit]
\lean{Complexity.cont_parse_double}
\label{Complexity.cont_parse_double}
\leanok
\uses{Complexity.choiceTapes, Complexity.contLoadCfg, Complexity.contPairTM, Complexity.cont_load_right, Turing.FinNDTM, Turing.MultiTapeTM.runFrom, Turing.MultiTapeTM.runFrom_succ_eq_step, Turing.MultiTapeTM.runFrom_zero}
\difficulty{5}
Let $N$ be a nondeterministic machine and suppose the physical input has the form $y = p\,b\,b\,r$ for a bit $b$. From the initial parser state at position $|p|+1$ with input buffer $x$, exactly two steps of the pair-loader machine append $b$ to the input buffer and move to position $|p|+3$, leaving the choice buffer empty and the output silent.
\end{lemma}

\begin{lemma}[Two-step recognition of the separator]
\lean{Complexity.cont_parse_separator}
\label{Complexity.cont_parse_separator}
\leanok
\uses{Complexity.choiceTapes, Complexity.contLoadCfg, Complexity.contPairTM, Complexity.cont_load_right, Turing.FinNDTM, Turing.MultiTapeTM.runFrom, Turing.MultiTapeTM.runFrom_succ_eq_step, Turing.MultiTapeTM.runFrom_zero}
\difficulty{4}
Let $N$ be a nondeterministic machine and suppose the physical input has the form $y = p\,0\,1\,u$. From the initial parser state at position $|p|+1$ with input buffer $x$, exactly two steps of the pair-loader machine enter the suffix-copy state at position $|p|+3$ without modifying either buffer or emitting output.
\end{lemma}

\begin{lemma}[Exact parse of the doubled prefix]
\lean{Complexity.cont_parse_run}
\label{Complexity.cont_parse_run}
\leanok
\uses{Complexity.contLoadCfg, Complexity.contPairTM, Complexity.contSplitWord, Complexity.cont_parse_double, Complexity.cont_parse_separator, Turing.FinNDTM, Turing.MultiTapeTM.runFrom, Turing.MultiTapeTM.runFrom_add, Turing.pairEncode}
\difficulty{5}
Let $N$ be a nondeterministic machine and suppose the physical input has the form $y = p\,e$ where $e$ is the self-delimiting encoding of a pair $(x, u)$. From the initial parser state at position $|p|+1$ with input buffer $a$, exactly $2|x| + 2$ steps of the pair-loader machine decode the doubled prefix and the separator, reaching the suffix-copy state at position $|p| + 2|x| + 3$ with input buffer $a\,x$ and the choice buffer still empty.
\end{lemma}

\begin{lemma}[Whole-machine contract on encoded pairs]
\lean{Complexity.cont_pair_computes}
\label{Complexity.cont_pair_computes}
\leanok
\uses{Complexity.choiceTapes, Complexity.contLoadCfg, Complexity.contPairTM, Complexity.cont_parse_run, Complexity.cont_suffix_run, Turing.Cfg.init, Turing.FinNDTM, Turing.FinTM.ComputesInTime, Turing.FinTM.ComputesInTime.mono, Turing.FinTM.computesInTime_iff, Turing.MultiTapeTM.initCfg, Turing.MultiTapeTM.runFrom_add, Turing.NDTM.initCfg, Turing.NDTM.runWith, Turing.pairEncode}
\difficulty{5}
Let $N$ be a nondeterministic Turing machine over the binary alphabet and let $x$ and $u$ be binary strings. On the self-delimiting encoding of the pair $(x,u)$, the pair-loader machine of $N$ halts from blank tapes within $3(\ell + 1)$ steps, where $\ell$ is the length of the encoded pair, and outputs the single bit recording whether the run of $N$ on $x$ under choice word $u$ has halted with output the one-bit string $1$.
\end{lemma}

\begin{lemma}[One-step rejection of the empty word]
\lean{Complexity.cont_pair_empty}
\label{Complexity.cont_pair_empty}
\leanok
\uses{Complexity.contPairTM, Turing.Cfg.init, Turing.Cfg.inputSymbol, Turing.FinNDTM, Turing.FinTM.ComputesInTime, Turing.FinTM.computesInTime_iff, Turing.MultiTapeTM.initCfg, Turing.MultiTapeTM.runFrom_succ_eq_step, Turing.MultiTapeTM.step}
\difficulty{2}
For every nondeterministic Turing machine $N$ over the binary alphabet, the pair-loader machine of $N$ on the empty input halts within one step, and its output is the single bit $0$.
\end{lemma}

\begin{definition}[Split emitter for the polynomial length equation]
\lean{Complexity.contSplitWord}
\label{Complexity.contSplitWord}
\leanok
\uses{Turing.pairEncode, Turing.solveSplit}
For natural numbers $C, c$ and a binary string $y$, the emitted word is the self-delimiting encoding of the pair formed by the first $i$ bits of $y$ and the remaining bits of $y$, when the split search solves $i + C\,(i+1)^c = |y|$; it is the empty word when the search fails.
\end{definition}

\begin{lemma}[Linear length bound on the emitted split]
\lean{Complexity.cont_split_length}
\label{Complexity.cont_split_length}
\leanok
\uses{Complexity.certificateSplit_spec, Complexity.contSplitWord, Complexity.cont_split_bridge, Turing.pairEncode, Turing.solveSplit}
\difficulty{3}
For all natural numbers $C, c$ and every binary string $y$, the word emitted by the polynomial split emitter on $y$ has length at most $2|y| + 2$.
\end{lemma}

\begin{lemma}[The loader decides the verifier on emitted splits]
\lean{Complexity.cont_split_answer}
\label{Complexity.cont_split_answer}
\leanok
\uses{Complexity.certificateSplit_spec, Complexity.choiceVerifier, Complexity.choiceVerifier_append, Complexity.choiceVerifier_no_split, Complexity.contPairTM, Complexity.contSplitWord, Complexity.cont_pair_computes, Complexity.cont_pair_empty, Complexity.cont_split_bridge, Complexity.cont_split_length, Turing.FinNDTM, Turing.FinTM.ComputesInTime, Turing.FinTM.ComputesInTime.mono, Turing.MultiTapeTM.indicator, Turing.NDTM.initCfg, Turing.NDTM.runWith, Turing.solveSplit}
\difficulty{5}
Let $N$ be a nondeterministic Turing machine over the binary alphabet and let $C, c$ be natural numbers. For every binary string $y$, the pair-loader machine of $N$, run on the word emitted by the polynomial split emitter for $y$, halts within $3(2|y|+3)$ steps, and its output is the indicator bit of membership of $y$ in the choice-word verifier language of $N$ with parameters $C, c$.
\end{lemma}

\begin{lemma}[The choice-word verifier is polynomial time]
\lean{Complexity.cont_choiceVerifier_mem_P}
\label{Complexity.cont_choiceVerifier_mem_P}
\leanok
\uses{Complexity.P, Complexity.choiceVerifier, Complexity.contPairTM, Complexity.contSplitWord, Complexity.cont_split_answer, Complexity.cont_split_length, Complexity.mem_P_iff, Turing.Cfg.init, Turing.FinNDTM, Turing.FinTM.ComputesInTime, Turing.FinTM.ComputesInTime.mono, Turing.FinTM.VirtualTag, Turing.FinTM.bufferedCompTM, Turing.FinTM.bufferedComp_start, Turing.FinTM.bufferedSecondCfg, Turing.FinTM.bufferedSecondCfg_run, Turing.FinTM.computesFunInTime_splitSolve, Turing.FinTM.computesInTime_iff, Turing.MultiTapeTM.indicator, Turing.MultiTapeTM.initCfg, Turing.MultiTapeTM.runFrom_add}
\difficulty{5}
For every nondeterministic Turing machine $N$ over the binary alphabet and all natural numbers $C$ and $c$, the choice-word verifier language of $N$ with parameters $C, c$ belongs to the class $\mathrm{P}$.
\end{lemma}

\begin{theorem}[Fixed-degree nondeterministic time lies in NP]
\lean{Complexity.ntime_poly_subset_NP}
\label{Complexity.ntime_poly_subset_NP}
\leanok
\uses{Complexity.NP, Complexity.NTIME, Complexity.choiceVerifier, Complexity.choice_certificate_iff, Complexity.cont_choiceVerifier_mem_P}
\difficulty{2}
For every natural number $c$, the nondeterministic time class $\mathrm{NTIME}(n^c + 1)$ is contained in $\mathrm{NP}$: every language decided by some nondeterministic Turing machine within time $a\,(n^c+1)$, for a constant $a$, admits a polynomial-time verifier with certificates of exact length $2a\,(n+1)^c$.
\end{theorem}

\subsection{The scheduler-driven guessing phase}

\begin{definition}[Selection of guess bits along a mask]
\lean{Complexity.contSelect}
\label{Complexity.contSelect}
\leanok
Given a mask and a word of physical choice bits, processed position by position, the selected word keeps the choice bit at every position where the mask is $1$ and drops it where the mask is $0$; positions beyond the shorter of the two lists contribute nothing.
\end{definition}

\begin{lemma}[Length of the selected word]
\lean{Complexity.cont_select_length}
\label{Complexity.cont_select_length}
\leanok
\uses{Complexity.contSelect}
\difficulty{3}
For every mask $m$ and every choice word $w$ with $|w| = |m|$, the word selected from $w$ along $m$ has length equal to the number of $1$ entries of $m$.
\end{lemma}

\begin{lemma}[Every word of the scheduled length is selected]
\lean{Complexity.cont_select_surjective}
\label{Complexity.cont_select_surjective}
\leanok
\uses{Complexity.contSelect}
\difficulty{4}
Let $m$ be a mask and let $u$ be a binary string whose length equals the number of $1$ entries of $m$. Then there is a choice word $w$ with $|w| = |m|$ whose selection along $m$ is exactly $u$.
\end{lemma}

\begin{definition}[Emission schedule of a deterministic run]
\lean{Complexity.contEmissionMask}
\label{Complexity.contEmissionMask}
\leanok
\uses{Turing.Cfg, Turing.FinTM, Turing.MultiTapeTM.outputSymbol, Turing.MultiTapeTM.step}
For a deterministic machine $M$, a configuration $d$ and a step count $t$, the emission mask is the length-$t$ bit list whose $i$-th entry records whether the machine emits an output symbol on the $(i+1)$-st step of its run from $d$.
\end{definition}

\begin{lemma}[One mask entry per transition]
\lean{Complexity.cont_mask_length}
\label{Complexity.cont_mask_length}
\leanok
\uses{Complexity.contEmissionMask, Turing.Cfg, Turing.FinTM}
\difficulty{2}
For every deterministic machine $M$, configuration $d$ and step count $t$, the emission mask of the run of $M$ from $d$ over $t$ steps has length exactly $t$.
\end{lemma}

\begin{lemma}[Marked positions count the emissions]
\lean{Complexity.cont_mask_count}
\label{Complexity.cont_mask_count}
\leanok
\uses{Complexity.contEmissionMask, Complexity.contGuessTM, Turing.Cfg, Turing.FinTM, Turing.MultiTapeTM.outputSymbol, Turing.MultiTapeTM.runFrom, Turing.MultiTapeTM.runFrom_succ_eq_step, Turing.MultiTapeTM.step, Turing.MultiTapeTM.step_output}
\difficulty{4}
Let $M$ be a deterministic machine, $d$ a configuration and $t$ a step count. The output length of $d$ plus the number of $1$ entries in the emission mask over $t$ steps equals the output length of the configuration reached after $t$ steps of $M$ from $d$.
\end{lemma}

\begin{definition}[Scheduler-driven nondeterministic guessing phase]
\lean{Complexity.contGuessTM}
\label{Complexity.contGuessTM}
\leanok
\uses{Turing.FinNDTM, Turing.FinTM, Turing.FinTM.controlAction, Turing.captureAction}
For a deterministic machine $M$ over the binary alphabet, a nondeterministic machine with one extra tape is defined: it runs $M$ step for step, and whenever $M$ emits an output symbol it instead appends the current nondeterministic choice bit to the extra capture tape; when $M$ halts, the phase stays in a live return state and makes no further changes.
\end{definition}

\begin{definition}[Invariant configuration of the guessing phase]
\lean{Complexity.contGuessCfg}
\label{Complexity.contGuessCfg}
\leanok
\uses{Complexity.contGuessTM, Turing.Cfg, Turing.FinTM, Turing.FinTM.bufferTape}
Given a configuration $d$ of a deterministic machine $M$ and a captured word $u$, the corresponding configuration of the guessing-phase machine of $M$ carries the state, input head and work tapes of $d$ unchanged, holds $u$ on the extra capture tape with its head at the right blank, and has empty physical output.
\end{definition}

\begin{lemma}[One guessing transition captures at most one bit]
\lean{Complexity.cont_guess_step}
\label{Complexity.cont_guess_step}
\leanok
\uses{Complexity.contGuessCfg, Complexity.contGuessTM, Turing.Action.apply, Turing.Cfg, Turing.Cfg.inputSymbol, Turing.Cfg.workTapeSymbols, Turing.FinTM, Turing.FinTM.bufferTape_append, Turing.FinTM.controlAction, Turing.FinTM.controlAction_apply, Turing.MultiTapeTM.outputSymbol, Turing.MultiTapeTM.step, Turing.MultiTapeTM.step_of_halt, Turing.NDTM.stepWith, Turing.captureAction, Turing.moveInputPos_zero}
\difficulty{5}
Let $M$ be a deterministic machine over the binary alphabet, $d$ one of its configurations, $u$ a captured word and $b$ a choice bit. One nondeterministic step of the guessing-phase machine of $M$ under $b$ maps the invariant configuration for $(d, u)$ to the invariant configuration for the stepped configuration of $M$, with $b$ appended to the captured word exactly when $M$ emits an output symbol at $d$.
\end{lemma}

\begin{lemma}[The guessing run realizes the emission mask]
\lean{Complexity.cont_guess_run}
\label{Complexity.cont_guess_run}
\leanok
\uses{Complexity.contEmissionMask, Complexity.contGuessCfg, Complexity.contGuessTM, Complexity.contSelect, Complexity.cont_guess_step, Turing.Cfg, Turing.FinTM, Turing.MultiTapeTM.runFrom, Turing.MultiTapeTM.runFrom_succ_eq_step, Turing.NDTM.runWith, Turing.NDTM.runWith_cons}
\difficulty{4}
Let $M$ be a deterministic machine over the binary alphabet, $d$ a configuration, $u$ a captured word and $w$ a word of choice bits. Running the guessing-phase machine of $M$ under $w$ from the invariant configuration for $(d,u)$ yields the invariant configuration for the $|w|$-step run of $M$ from $d$, with the selection of $w$ along the emission mask appended to the captured word.
\end{lemma}

\begin{lemma}[Initial configuration of the guessing phase]
\lean{Complexity.cont_guess_initial}
\label{Complexity.cont_guess_initial}
\leanok
\uses{Complexity.contGuessCfg, Complexity.contGuessTM, Turing.Cfg.init, Turing.FinTM, Turing.MultiTapeTM.initCfg, Turing.NDTM.initCfg}
\difficulty{3}
For every deterministic machine $M$ over the binary alphabet and every input $x$, the blank-tape initial configuration of the guessing-phase machine of $M$ on $x$ equals the invariant configuration for the initial configuration of $M$ on $x$ with empty captured word.
\end{lemma}

\begin{lemma}[Extraction and coverage of guessed words]
\lean{Complexity.cont_guess_coverage}
\label{Complexity.cont_guess_coverage}
\leanok
\uses{Complexity.contEmissionMask, Complexity.contGuessCfg, Complexity.contGuessTM, Complexity.contSelect, Complexity.cont_guess_initial, Complexity.cont_guess_run, Complexity.cont_mask_count, Complexity.cont_mask_length, Complexity.cont_select_length, Complexity.cont_select_surjective, Turing.Cfg.init, Turing.FinTM, Turing.FinTM.ComputesInTime, Turing.FinTM.computesInTime_iff, Turing.MultiTapeTM.initCfg, Turing.MultiTapeTM.runFrom, Turing.NDTM.initCfg, Turing.NDTM.runWith}
\difficulty{5}
Let $M$ be a deterministic machine over the binary alphabet that computes an output $v$ on input $x$ within $T$ steps. Then every choice word of length $T$ drives the guessing-phase machine of $M$ from its initial configuration on $x$ to the invariant configuration for the $T$-step run of $M$, with some captured word of length $|v|$; and conversely, every word $u$ with $|u| = |v|$ is the captured word of that configuration for some choice word of length $T$.
\end{lemma}

\begin{lemma}[Normalizing the compiled polynomial envelope]
\lean{Complexity.cont_guess_time_bound}
\label{Complexity.cont_guess_time_bound}
\leanok
\uses{Complexity.succ_pow_le}
\difficulty{4}
For all natural numbers $C, c, K, r$ and $n$, the bound $K\,(n + C(n+1)^c + 1)^r \le K\,(C+1)^r\, 2^{\,r \cdot \max(1,c)} \cdot (n^{\,r \cdot \max(1,c)} + 1)$ holds.
\end{lemma}

\begin{lemma}[Landing a compiled decider in an NTIME component]
\lean{Complexity.cont_guess_normalize}
\label{Complexity.cont_guess_normalize}
\leanok
\uses{Complexity.NTIME, Complexity.acceptsWithin_iff_of_halts, Complexity.cont_guess_time_bound, Turing.FinNDTM, Turing.FinNDTM.DecidesInTime, Turing.NDTM.HaltsWithin.mono}
\difficulty{3}
Let $L$ be a language over the binary alphabet and $C, c$ natural numbers. If some nondeterministic machine decides $L$ within time $K\,(n + C(n+1)^c + 1)^r$ for some constants $K$ and $r$, then $L$ belongs to $\bigcup_{e} \mathrm{NTIME}(n^e + 1)$.
\end{lemma}

\begin{lemma}[A concrete guessing phase for each polynomial length]
\lean{Complexity.cont_poly_guess_phase}
\label{Complexity.cont_poly_guess_phase}
\leanok
\uses{Complexity.contGuessCfg, Complexity.contGuessTM, Complexity.cont_guess_coverage, Turing.FinTM, Turing.FinTM.computesFunInTime_polyUnary, Turing.MultiTapeTM.initCfg, Turing.MultiTapeTM.runFrom, Turing.NDTM.initCfg, Turing.NDTM.runWith}
\difficulty{3}
For all natural numbers $C$ and $c$ there exist a deterministic scheduler machine $M$ and a constant $A$ such that for every $n$, on the unary input of length $n$ and with $T = A(n+1)^{c+1}$: every choice word of length $T$ drives the guessing-phase machine of $M$ to the completed $T$-step configuration while capturing some word of length exactly $C(n+1)^c$, and every word of that length is captured by some choice word of length $T$.
\end{lemma}

\subsection{Unary normalization of schedulers}

\begin{definition}[Input-read normalization of a machine]
\lean{Complexity.b2UnaryTM}
\label{Complexity.b2UnaryTM}
\leanok
\uses{Turing.FinTM}
For a deterministic machine $M$ over the binary alphabet, a machine with the same states and tapes is defined whose transition table first replaces every non-blank input symbol it reads by the bit $1$ and otherwise acts exactly as $M$; blanks are preserved, so only interior input reads are normalized.
\end{definition}

\begin{definition}[Viewing a configuration over the unary input]
\lean{Complexity.b2UnaryCfg}
\label{Complexity.b2UnaryCfg}
\leanok
\uses{Turing.Cfg}
A configuration on an input $x$ is reinterpreted as a configuration on the all-ones word of length $|x|$, keeping its state, input head position, work tapes, work heads and output unchanged.
\end{definition}

\begin{lemma}[Reads after unary reinterpretation]
\lean{Complexity.b2_unary_read}
\label{Complexity.b2_unary_read}
\leanok
\uses{Complexity.b2UnaryCfg, Turing.Cfg, Turing.Cfg.inputSymbol}
\difficulty{3}
For every configuration $d$ on an input $x$, the input symbol read by the reinterpretation of $d$ over the all-ones word of length $|x|$ equals the symbol read by $d$ with every non-blank bit replaced by $1$.
\end{lemma}

\begin{lemma}[Actions commute with unary reinterpretation]
\lean{Complexity.b2_unary_apply}
\label{Complexity.b2_unary_apply}
\leanok
\uses{Complexity.b2UnaryCfg, Turing.Action, Turing.Action.apply, Turing.Cfg, Turing.moveInputPos}
\difficulty{3}
For every tape action $a$ and every configuration $d$ on an input $x$, reinterpreting the result of applying $a$ to $d$ over the all-ones word equals applying $a$ to the reinterpretation of $d$; input-head clamping depends only on the input length.
\end{lemma}

\begin{lemma}[One normalized step equals one unary step]
\lean{Complexity.b2_unary_step}
\label{Complexity.b2_unary_step}
\leanok
\uses{Complexity.b2UnaryCfg, Complexity.b2UnaryTM, Complexity.b2_unary_apply, Complexity.b2_unary_read, Turing.Action.apply, Turing.Cfg, Turing.Cfg.inputSymbol, Turing.Cfg.workTapeSymbols, Turing.FinTM, Turing.MultiTapeTM.step}
\difficulty{4}
Let $M$ be a deterministic machine over the binary alphabet and $d$ a configuration on an input $x$. The unary reinterpretation of one step of the input-normalized variant of $M$ at $d$ equals one native step of $M$ at the unary reinterpretation of $d$, including the absorbing halted case.
\end{lemma}

\begin{lemma}[Exact unary lockstep at every time]
\lean{Complexity.b2_unary_run}
\label{Complexity.b2_unary_run}
\leanok
\uses{Complexity.b2UnaryCfg, Complexity.b2UnaryTM, Complexity.b2_unary_step, Turing.Cfg, Turing.FinTM, Turing.MultiTapeTM.runFrom, Turing.MultiTapeTM.runFrom_succ_eq_step}
\difficulty{3}
Let $M$ be a deterministic machine over the binary alphabet, $d$ a configuration on an input $x$, and $t$ a step count. The unary reinterpretation of the $t$-step run of the input-normalized variant of $M$ from $d$ equals the $t$-step run of $M$ from the reinterpretation of $d$.
\end{lemma}

\begin{lemma}[Unary reinterpretation of the initial configuration]
\lean{Complexity.b2_unary_initial}
\label{Complexity.b2_unary_initial}
\leanok
\uses{Complexity.b2UnaryCfg, Complexity.b2UnaryTM, Turing.Cfg.init, Turing.FinTM, Turing.MultiTapeTM.initCfg}
\difficulty{2}
For every deterministic machine $M$ over the binary alphabet and every input $x$, the reinterpretation of the initial configuration of the input-normalized variant of $M$ on $x$ equals the initial configuration of $M$ on the all-ones word of length $|x|$.
\end{lemma}

\begin{lemma}[Emission schedule survives normalization]
\lean{Complexity.b2_unary_mask}
\label{Complexity.b2_unary_mask}
\leanok
\uses{Complexity.b2UnaryCfg, Complexity.b2UnaryTM, Complexity.b2_unary_read, Complexity.b2_unary_step, Complexity.contEmissionMask, Turing.Cfg, Turing.FinTM, Turing.MultiTapeTM.outputSymbol}
\difficulty{4}
Let $M$ be a deterministic machine over the binary alphabet, $d$ a configuration and $t$ a step count. The emission mask of the input-normalized variant of $M$ over $t$ steps from $d$ equals the emission mask of $M$ over $t$ steps from the unary reinterpretation of $d$.
\end{lemma}

\begin{lemma}[Transfer of unary computations to all inputs]
\lean{Complexity.b2_unary_computes}
\label{Complexity.b2_unary_computes}
\leanok
\uses{Complexity.b2UnaryTM, Complexity.b2_unary_initial, Complexity.b2_unary_run, Turing.FinTM, Turing.FinTM.ComputesInTime, Turing.FinTM.computesInTime_iff, Turing.MultiTapeTM.initCfg}
\difficulty{3}
Let $M$ be a deterministic machine over the binary alphabet and suppose $M$ computes an output $v$ within $t$ steps on the all-ones input of length $|x|$. Then the input-normalized variant of $M$ computes $v$ within $t$ steps on the input $x$ itself.
\end{lemma}

\begin{lemma}[Length-indexed first halting times]
\lean{Complexity.b2_unary_first}
\label{Complexity.b2_unary_first}
\leanok
\uses{Complexity.b2UnaryTM, Complexity.b2_unary_computes, Complexity.b2_unary_initial, Complexity.b2_unary_run, Turing.FinTM, Turing.FinTM.ComputesInTime, Turing.FinTM.computesInTime_iff, Turing.MultiTapeTM.initCfg, Turing.MultiTapeTM.runFrom, Turing.MultiTapeTM.runFrom_output_eq_of_halt}
\difficulty{5}
Let $M$ be a deterministic machine over the binary alphabet such that for every $n$, on the all-ones input of length $n$, $M$ computes a word $f(n)$ within $T(n)$ steps. Then there is a function $\tau$ with $\tau(n) \le T(n)$ for all $n$ such that on every input $x$, the input-normalized variant of $M$ computes $f(|x|)$ within $\tau(|x|)$ steps and is still live at every time before $\tau(|x|)$.
\end{lemma}

\subsection{The reverse host and Theorem 2.6}

\begin{definition}[Four disjoint tape banks of the reverse host]
\lean{Complexity.b2Slots}
\label{Complexity.b2Slots}
\leanok
\uses{Turing.FinTM}
Given deterministic machines $S$ and $V$, values for the work tapes of $S$, an assembly value, values for the work tapes of $V$ and a capture value, the combined assignment is the function that places the scheduler bank first with the assembly buffer as its final slot, followed by the verifier bank with the capture tape as its final slot.
\end{definition}

\begin{definition}[The reverse-host nondeterministic machine]
\lean{Complexity.b2Host}
\label{Complexity.b2Host}
\leanok
\uses{Complexity.b2Slots, Complexity.b2UnaryTM, Complexity.captureEmission, Complexity.contGuessTM, Turing.FinNDTM, Turing.FinTM, Turing.FinTM.controlAction, Turing.FinTM.leftAction, Turing.FinTM.virtualMove, Turing.FinTM.virtualNextTag}
For deterministic machines $S$ (a scheduler) and $V$ (a verifier) over the binary alphabet, a nondeterministic machine is defined that copies its input to an assembly buffer and rewinds; runs the input-normalized scheduler, appending one nondeterministic choice bit to the assembly buffer at each scheduler emission; rewinds the assembly buffer; then runs $V$ with its input relocated to the assembled word under a guarded virtual head, capturing every emission of $V$; and finally halts with the single bit recording whether the captured output is exactly the one-bit string $1$. Outside the live guessing states its two transition tables coincide.
\end{definition}

\begin{definition}[Guessing-phase configuration of the host]
\lean{Complexity.b2GuessCfg}
\label{Complexity.b2GuessCfg}
\leanok
\uses{Complexity.b2Host, Complexity.b2UnaryTM, Complexity.contGuessCfg, Turing.Cfg, Turing.FinTM, Turing.FinTM.leftCfg}
Given a scheduler configuration $d$ on an input $x$ and a guessed word $u$, the corresponding configuration is the host configuration that embeds the guessing-phase invariant for $d$ with captured word $x\,u$ in the scheduler bank and keeps the entire verifier bank blank.
\end{definition}

\begin{lemma}[Choice independence outside guessing states]
\lean{Complexity.b2_tables_coincide}
\label{Complexity.b2_tables_coincide}
\leanok
\uses{Complexity.b2Host, Turing.FinTM}
\difficulty{2}
Let $S$ and $V$ be deterministic machines and let $q$ be any control state of the reverse host built from $S$ and $V$ that is not a live guessing state. For every input symbol and every tuple of work-tape symbols, the two transition tables of the host agree at $q$.
\end{lemma}

\begin{lemma}[A live host guessing step]
\lean{Complexity.b2_guess_step}
\label{Complexity.b2_guess_step}
\leanok
\uses{Complexity.b2GuessCfg, Complexity.b2Host, Complexity.b2UnaryTM, Complexity.contGuessCfg, Complexity.contGuessTM, Complexity.cont_guess_step, Turing.Action.apply, Turing.Cfg, Turing.Cfg.inputSymbol, Turing.Cfg.workTapeSymbols, Turing.FinTM, Turing.FinTM.leftAction, Turing.FinTM.leftCfg, Turing.FinTM.leftCfg_apply, Turing.MultiTapeTM.outputSymbol, Turing.MultiTapeTM.step, Turing.NDTM.stepWith}
\difficulty{5}
Let $S, V$ be deterministic machines, $d$ a non-halted scheduler configuration on an input $x$, $u$ a guessed word and $b$ a choice bit. One step of the reverse host under $b$ maps the guessing-phase configuration for $(d,u)$ to the guessing-phase configuration for the stepped configuration of the input-normalized variant of $S$, with $b$ appended to the guessed word exactly when that machine emits at $d$.
\end{lemma}

\begin{lemma}[Host guessing run under a liveness guard]
\lean{Complexity.b2_guess_run}
\label{Complexity.b2_guess_run}
\leanok
\uses{Complexity.b2GuessCfg, Complexity.b2Host, Complexity.b2UnaryTM, Complexity.b2_guess_step, Complexity.contEmissionMask, Complexity.contSelect, Turing.Cfg, Turing.FinTM, Turing.MultiTapeTM.runFrom, Turing.MultiTapeTM.runFrom_succ_eq_step, Turing.MultiTapeTM.step, Turing.NDTM.runWith, Turing.NDTM.runWith_cons}
\difficulty{4}
Let $S, V$ be deterministic machines, $d$ a scheduler configuration on an input $x$, $u$ a guessed word, and $w$ a word of choice bits such that the input-normalized scheduler is still live at every time before $|w|$. Running the reverse host under $w$ from the guessing-phase configuration for $(d,u)$ gives the guessing-phase configuration for the $|w|$-step scheduler run, with the selection of $w$ along the scheduler emission mask appended to the guessed word.
\end{lemma}

\begin{definition}[Loader configuration of the reverse host]
\lean{Complexity.b2LoadCfg}
\label{Complexity.b2LoadCfg}
\leanok
\uses{Complexity.b2Host, Complexity.b2Slots, Turing.Cfg, Turing.FinTM, Turing.FinTM.bufferTape}
Given a loader state index, a physical input position, a copied prefix and an assembly head value, the corresponding host configuration has every scheduler and verifier tape blank, the copied prefix on the assembly buffer, and empty physical output; the physical input is retained verbatim.
\end{definition}

\begin{lemma}[One copy transition of the host]
\lean{Complexity.b2_copy_step}
\label{Complexity.b2_copy_step}
\leanok
\uses{Complexity.b2Host, Complexity.b2LoadCfg, Complexity.b2Slots, Turing.Action.apply, Turing.Cfg.inputSymbol, Turing.FinTM, Turing.FinTM.bufferTape_append, Turing.NDTM.stepWith, Turing.inputSymbolInner, Turing.moveInputPos, Turing.moveInputPos_pos_of_ne_right}
\difficulty{4}
Let $S, V$ be deterministic machines and $x$ an input with position index $i < |x|$. For every choice bit, one step of the reverse host from the copy-state loader configuration holding the first $i$ bits of $x$ appends the bit of $x$ at index $i$ to the assembly buffer and advances the input head by one.
\end{lemma}

\begin{lemma}[Choice-independent copying of a prefix]
\lean{Complexity.b2_copy_run}
\label{Complexity.b2_copy_run}
\leanok
\uses{Complexity.b2Host, Complexity.b2LoadCfg, Complexity.b2_copy_step, Turing.FinTM, Turing.NDTM.runWith, Turing.NDTM.runWith_cons}
\difficulty{4}
Let $S, V$ be deterministic machines, $x$ an input, and $w$ any word of choice bits with $i + |w| \le |x|$. Running the reverse host under $w$ from the copy-state loader configuration holding the first $i$ bits of $x$ yields the copy-state configuration holding the first $i + |w|$ bits, one bit copied per step regardless of the choices.
\end{lemma}

\begin{lemma}[Dispatch from the left input blank]
\lean{Complexity.b2_rewind_done}
\label{Complexity.b2_rewind_done}
\leanok
\uses{Complexity.b2GuessCfg, Complexity.b2Host, Complexity.b2LoadCfg, Complexity.b2Slots, Complexity.b2UnaryTM, Complexity.contGuessCfg, Turing.Action.apply, Turing.Cfg.init, Turing.Cfg.inputSymbol, Turing.FinTM, Turing.FinTM.controlAction_apply, Turing.FinTM.leftCfg, Turing.MultiTapeTM.initCfg, Turing.NDTM.stepWith, Turing.moveInputPos}
\difficulty{4}
Let $S, V$ be deterministic machines and $x$ an input. From the rewind-state loader configuration with input head at the left blank and assembly buffer exactly $x$, one host step (for any choice bit) yields the guessing-phase configuration for the initial configuration of the input-normalized scheduler on $x$ with empty guessed word.
\end{lemma}

\begin{lemma}[One left move of the input rewind]
\lean{Complexity.b2_rewind_step}
\label{Complexity.b2_rewind_step}
\leanok
\uses{Complexity.b2Host, Complexity.b2LoadCfg, Turing.Action.apply, Turing.Cfg.inputSymbol, Turing.FinTM, Turing.FinTM.controlAction_apply, Turing.FinTM.moveInputPos_neg_val, Turing.NDTM.stepWith, Turing.inputSymbolInner, Turing.moveInputPos}
\difficulty{3}
Let $S, V$ be deterministic machines and $x$ an input with $j < |x|$. For any choice bit, one host step from the rewind-state loader configuration with input head at position $j+1$ and assembly buffer $x$ moves the input head to position $j$, leaving all tapes unchanged.
\end{lemma}

\begin{lemma}[Exact cost of the physical input rewind]
\lean{Complexity.b2_rewind_run}
\label{Complexity.b2_rewind_run}
\leanok
\uses{Complexity.b2GuessCfg, Complexity.b2Host, Complexity.b2LoadCfg, Complexity.b2UnaryTM, Complexity.b2_rewind_done, Complexity.b2_rewind_step, Turing.FinTM, Turing.MultiTapeTM.initCfg, Turing.NDTM.runWith, Turing.NDTM.runWith_cons}
\difficulty{4}
Let $S, V$ be deterministic machines, $x$ an input, $j \le |x|$, and $w$ any word of choice bits with $|w| = j + 1$. Running the reverse host under $w$ from the rewind-state loader configuration with input head at position $j$ and assembly buffer $x$ reaches the guessing-phase configuration for the initial normalized-scheduler configuration on $x$ with empty guessed word.
\end{lemma}

\begin{lemma}[Initial configuration of the reverse host]
\lean{Complexity.b2_initial}
\label{Complexity.b2_initial}
\leanok
\uses{Complexity.b2Host, Complexity.b2LoadCfg, Complexity.b2Slots, Turing.FinTM, Turing.NDTM.initCfg}
\difficulty{3}
For all deterministic machines $S, V$ and every input $x$, the blank-tape initial configuration of the reverse host built from $S$ and $V$ equals the copy-state loader configuration with empty assembly buffer, input head at position one and assembly head at zero.
\end{lemma}

\begin{lemma}[Boundary turn after copying]
\lean{Complexity.b2_copy_done}
\label{Complexity.b2_copy_done}
\leanok
\uses{Complexity.b2Host, Complexity.b2LoadCfg, Turing.Action.apply, Turing.Cfg.inputSymbol, Turing.FinTM, Turing.FinTM.controlAction_apply, Turing.FinTM.inputSymbol_at, Turing.FinTM.moveInputPos_neg_val, Turing.NDTM.stepWith, Turing.moveInputPos}
\difficulty{3}
Let $S, V$ be deterministic machines and $x$ an input. For any choice bit, one host step from the copy-state loader configuration at the right input blank with assembly buffer $x$ yields the rewind-state configuration with the input head moved left to position $|x|$.
\end{lemma}

\begin{lemma}[Startup contract of the reverse host]
\lean{Complexity.b2_start}
\label{Complexity.b2_start}
\leanok
\uses{Complexity.b2GuessCfg, Complexity.b2Host, Complexity.b2LoadCfg, Complexity.b2UnaryTM, Complexity.b2_copy_done, Complexity.b2_copy_run, Complexity.b2_initial, Complexity.b2_rewind_run, Turing.FinTM, Turing.MultiTapeTM.initCfg, Turing.NDTM.initCfg, Turing.NDTM.runWith, Turing.NDTM.runWith_append, Turing.NDTM.runWith_cons}
\difficulty{5}
Let $S, V$ be deterministic machines, $x$ an input, and $w$ any word of choice bits with $|w| = 2|x| + 2$. Running the reverse host under $w$ from its blank-tape initial configuration on $x$ reaches the guessing-phase configuration for the initial configuration of the input-normalized scheduler on $x$ with empty guessed word; all startup choices are ignored.
\end{lemma}

\begin{lemma}[The tape contents determine the captured word]
\lean{Complexity.b2_guess_word_injective}
\label{Complexity.b2_guess_word_injective}
\leanok
\uses{Complexity.contGuessCfg, Turing.Cfg, Turing.FinTM}
\difficulty{3}
Let $M$ be a deterministic machine, $d$ one of its configurations, and $u, v$ captured words. If the guessing-phase invariant configurations for $(d,u)$ and $(d,v)$ are equal, then $u = v$.
\end{lemma}

\begin{lemma}[Host-level extraction and coverage of guesses]
\lean{Complexity.b2_guess_coverage}
\label{Complexity.b2_guess_coverage}
\leanok
\uses{Complexity.b2GuessCfg, Complexity.b2Host, Complexity.b2UnaryTM, Complexity.b2_guess_run, Complexity.b2_guess_word_injective, Complexity.contGuessCfg, Complexity.contGuessTM, Complexity.cont_guess_coverage, Complexity.cont_guess_initial, Complexity.cont_guess_run, Turing.FinTM, Turing.FinTM.ComputesInTime, Turing.MultiTapeTM.initCfg, Turing.MultiTapeTM.runFrom, Turing.NDTM.initCfg, Turing.NDTM.runWith}
\difficulty{5}
Let $S, V$ be deterministic machines and suppose the input-normalized scheduler on an input $x$ computes an output $v$ within $T$ steps while remaining live at all earlier times. Then every choice word of length $T$ drives the reverse host from the start-of-guessing configuration to the completed-scheduler guessing configuration with some guessed word of length $|v|$, and every word of length $|v|$ arises this way from some choice word of length $T$.
\end{lemma}

\begin{definition}[Verifier-phase configuration of the host]
\lean{Complexity.b2VerifyCfg}
\label{Complexity.b2VerifyCfg}
\leanok
\uses{Complexity.b2Host, Complexity.b2Slots, Complexity.capturedSummary, Turing.Cfg, Turing.FinTM, Turing.FinTM.bufferTape}
Given a verifier configuration $d$ on an assembled word $y$, a boundary tag, a physical input position and frozen scheduler tapes and heads, the corresponding configuration is the host configuration that stores the verifier state, tag and output summary in control, $y$ on the assembly buffer serving as virtual input, the verifier work tapes in their own bank, the complete emitted output on the capture tape, and empty physical output.
\end{definition}

\begin{definition}[Assembly-rewind configuration of the host]
\lean{Complexity.b2ReadyCfg}
\label{Complexity.b2ReadyCfg}
\leanok
\uses{Complexity.b2Host, Complexity.b2Slots, Turing.Cfg, Turing.FinTM, Turing.FinTM.bufferTape}
Given a completed scheduler configuration $d$ on an input $x$, an assembled word $y$ and an assembly head value, the corresponding host configuration is in the assembly-rewind control state, retains the scheduler bank of $d$ and the word $y$ on the assembly buffer, and keeps the verifier bank blank.
\end{definition}

\begin{lemma}[One relocated verifier step]
\lean{Complexity.b2_verify_step}
\label{Complexity.b2_verify_step}
\leanok
\uses{Complexity.b2Host, Complexity.b2Slots, Complexity.b2VerifyCfg, Complexity.captureEmission, Complexity.captureEmission_correct, Complexity.capturedSummary, Turing.Action.apply, Turing.Cfg, Turing.Cfg.inputSymbol, Turing.Cfg.workTapeSymbols, Turing.FinTM, Turing.FinTM.VirtualTag, Turing.FinTM.bufferTape_append, Turing.FinTM.bufferTape_inputSymbol, Turing.FinTM.virtualMove, Turing.FinTM.virtualMove_correct, Turing.FinTM.virtualNextTag, Turing.MultiTapeTM.step, Turing.NDTM.stepWith, Turing.moveInputPos_zero}
\difficulty{6}
Let $S, V$ be deterministic machines, $d$ a non-halted verifier configuration on an assembled word $y$ carrying a valid virtual boundary tag, and fix frozen scheduler data and a physical position. For every choice bit, one host step maps the verifier-phase configuration for $d$ to the verifier-phase configuration for the stepped verifier, with some valid updated tag; reads come from the assembly buffer and every emission of $V$ is captured, including one on a halting transition.
\end{lemma}

\begin{lemma}[Relocated verifier run under a liveness guard]
\lean{Complexity.b2_verify_run}
\label{Complexity.b2_verify_run}
\leanok
\uses{Complexity.b2Host, Complexity.b2VerifyCfg, Complexity.b2_verify_step, Turing.Cfg, Turing.FinTM, Turing.FinTM.VirtualTag, Turing.MultiTapeTM.runFrom, Turing.MultiTapeTM.runFrom_succ_eq_step, Turing.MultiTapeTM.step, Turing.NDTM.runWith, Turing.NDTM.runWith_cons}
\difficulty{4}
Let $S, V$ be deterministic machines, $d$ a verifier configuration with a valid boundary tag, and $w$ a choice word such that the verifier is live at every time before $|w|$. Running the host under $w$ from the verifier-phase configuration for $d$ yields the verifier-phase configuration for the $|w|$-step verifier run, with a valid tag; the choices are ignored.
\end{lemma}

\begin{lemma}[Verdict transition after the verifier halts]
\lean{Complexity.b2_verify_finish}
\label{Complexity.b2_verify_finish}
\leanok
\uses{Complexity.b2Host, Complexity.b2VerifyCfg, Complexity.capturedSummary_true, Turing.Cfg, Turing.FinTM, Turing.NDTM.stepWith}
\difficulty{2}
Let $S, V$ be deterministic machines and $d$ a halted verifier configuration. For any choice bit, one host step from the verifier-phase configuration for $d$ halts the host and emits the single bit recording whether the complete output of $d$ is exactly the one-bit string $1$.
\end{lemma}

\begin{lemma}[Dispatch from the scheduler-return state]
\lean{Complexity.b2_guess_return}
\label{Complexity.b2_guess_return}
\leanok
\uses{Complexity.b2GuessCfg, Complexity.b2Host, Complexity.b2ReadyCfg, Complexity.b2Slots, Complexity.b2UnaryTM, Complexity.contGuessCfg, Turing.Cfg, Turing.FinTM, Turing.FinTM.leftCfg, Turing.NDTM.stepWith, Turing.moveInputPos_zero}
\difficulty{4}
Let $S, V$ be deterministic machines, $d$ a halted scheduler configuration on an input $x$, and $u$ a guessed word. For any choice bit, one host step from the guessing-phase configuration for $(d,u)$ yields the assembly-rewind configuration holding the assembled word $x\,u$ with assembly head at its last cell.
\end{lemma}

\begin{lemma}[One left move of the assembly rewind]
\lean{Complexity.b2_ready_step}
\label{Complexity.b2_ready_step}
\leanok
\uses{Complexity.b2Host, Complexity.b2ReadyCfg, Complexity.b2Slots, Turing.Action.apply, Turing.Cfg, Turing.Cfg.workTapeSymbols, Turing.FinTM, Turing.NDTM.stepWith, Turing.moveInputPos_zero}
\difficulty{4}
Let $S, V$ be deterministic machines, $d$ a scheduler configuration, $y$ an assembled word and $j < |y|$. For any choice bit, one host step from the assembly-rewind configuration with head at cell $j$ moves the head to cell $j - 1$, preserving $y$ and the scheduler bank.
\end{lemma}

\begin{lemma}[Verifier startup from the assembly blank]
\lean{Complexity.b2_ready_done}
\label{Complexity.b2_ready_done}
\leanok
\uses{Complexity.b2Host, Complexity.b2ReadyCfg, Complexity.b2Slots, Complexity.b2VerifyCfg, Turing.Action.apply, Turing.Cfg, Turing.Cfg.init, Turing.Cfg.workTapeSymbols, Turing.FinTM, Turing.MultiTapeTM.initCfg, Turing.NDTM.stepWith, Turing.moveInputPos_zero}
\difficulty{4}
Let $S, V$ be deterministic machines, $d$ a scheduler configuration and $y$ an assembled word. For any choice bit, one host step from the assembly-rewind configuration with head at cell $-1$ yields the verifier-phase configuration for the initial configuration of $V$ on $y$, with interior boundary tag, blank verifier tapes and empty captured output.
\end{lemma}

\begin{lemma}[Exact cost of the assembly rewind]
\lean{Complexity.b2_ready_run}
\label{Complexity.b2_ready_run}
\leanok
\uses{Complexity.b2Host, Complexity.b2ReadyCfg, Complexity.b2VerifyCfg, Complexity.b2_ready_done, Complexity.b2_ready_step, Turing.Cfg, Turing.FinTM, Turing.MultiTapeTM.initCfg, Turing.NDTM.runWith, Turing.NDTM.runWith_cons}
\difficulty{4}
Let $S, V$ be deterministic machines, $d$ a scheduler configuration, $y$ an assembled word, $j \le |y|$, and $w$ any choice word with $|w| = j + 1$. Running the host under $w$ from the assembly-rewind configuration with head at cell $j - 1$ reaches the verifier-phase configuration for the initial configuration of $V$ on $y$.
\end{lemma}

\begin{lemma}[From scheduler halt to relocated verifier start]
\lean{Complexity.b2_assembly}
\label{Complexity.b2_assembly}
\leanok
\uses{Complexity.b2GuessCfg, Complexity.b2Host, Complexity.b2VerifyCfg, Complexity.b2_guess_return, Complexity.b2_ready_run, Turing.Cfg, Turing.FinTM, Turing.MultiTapeTM.initCfg, Turing.NDTM.runWith, Turing.NDTM.runWith_cons}
\difficulty{3}
Let $S, V$ be deterministic machines, $d$ a halted scheduler configuration on an input $x$, $u$ a guessed word, and $w$ any choice word with $|w| = |x\,u| + 2$. Running the host under $w$ from the guessing-phase configuration for $(d,u)$ installs the verifier-phase configuration for the initial configuration of $V$ on the assembled word $x\,u$.
\end{lemma}

\begin{lemma}[All-branch timed completion of the verifier phase]
\lean{Complexity.b2_verify_timed}
\label{Complexity.b2_verify_timed}
\leanok
\uses{Complexity.b2Host, Complexity.b2VerifyCfg, Complexity.b2_verify_finish, Complexity.b2_verify_run, Turing.Cfg.init, Turing.FinTM, Turing.FinTM.ComputesInTime, Turing.FinTM.ComputesInTime.output_unique, Turing.FinTM.VirtualTag, Turing.FinTM.computesInTime_iff, Turing.MultiTapeTM.initCfg, Turing.MultiTapeTM.runFrom, Turing.NDTM.runWith, Turing.NDTM.runWith_append, Turing.NDTM.runWith_cons, Turing.NDTM.runWith_of_halt, Turing.NDTM.stepWith}
\difficulty{5}
Let $S, V$ be deterministic machines and suppose $V$ computes the single-bit output $b$ on a word $y$ within $T$ steps. For every choice word $w$ with $|w| = T + 1$, running the host under $w$ from the verifier-phase start on $y$ halts with output exactly the one-bit string $b$.
\end{lemma}

\begin{lemma}[Timed completion of assembly plus verification]
\lean{Complexity.b2_finish}
\label{Complexity.b2_finish}
\leanok
\uses{Complexity.b2GuessCfg, Complexity.b2Host, Complexity.b2_assembly, Complexity.b2_verify_timed, Turing.Cfg, Turing.FinTM, Turing.FinTM.ComputesInTime, Turing.NDTM.runWith, Turing.NDTM.runWith_append}
\difficulty{4}
Let $S, V$ be deterministic machines, $d$ a halted scheduler configuration on an input $x$, $u$ a guessed word, and suppose $V$ computes the single-bit output $b$ on $x\,u$ within $T$ steps. Then for every choice word $w$ with $|w| = |x\,u| + 2 + (T+1)$, running the host under $w$ from the guessing-phase configuration for $(d,u)$ halts with output exactly the one-bit string $b$.
\end{lemma}

\begin{lemma}[Integrated extraction and coverage of the host]
\lean{Complexity.b2_host_contract}
\label{Complexity.b2_host_contract}
\leanok
\uses{Complexity.b2GuessCfg, Complexity.b2Host, Complexity.b2UnaryTM, Complexity.b2_finish, Complexity.b2_guess_coverage, Complexity.b2_start, Turing.FinTM, Turing.FinTM.ComputesInTime, Turing.FinTM.DecidesInTime, Turing.FinTM.computesInTime_iff, Turing.MultiTapeTM.indicator, Turing.MultiTapeTM.initCfg, Turing.MultiTapeTM.runFrom, Turing.NDTM.initCfg, Turing.NDTM.runWith, Turing.NDTM.runWith_append}
\difficulty{6}
Let $S$ and $M$ be deterministic machines, $V$ a language, and suppose the input-normalized scheduler on an input $x$ computes a word $v$ within $T$ steps while live earlier, and that $M$ decides $V$ within time $A(m+1)^d$ on inputs of length $m$. Set $H = 2|x| + 2 + T + (|x| + |v| + 2 + (A\,(|x|+|v|+1)^d + 1))$. Then every choice word of length $H$ makes the reverse host built from $S$ and $M$ halt with output the indicator bit of $x\,u \in V$ for some guessed word $u$ of length $|v|$; and conversely, for every $u$ of length $|v|$ some choice word of length $H$ produces a halted run with output the indicator bit of $x\,u \in V$.
\end{lemma}

\begin{lemma}[Polynomial envelope of the host ledger]
\lean{Complexity.b2_host_bound}
\label{Complexity.b2_host_bound}
\leanok
\difficulty{4}
For all natural numbers $C, c, B, A, d, n$ and $T$ with $T \le B(n+1)^{c+1}$, the total $2n + 2 + T + (n + C(n+1)^c + 2 + (A\,(n + C(n+1)^c + 1)^d + 1))$ is at most $(B + A + 5)\,(n + C(n+1)^c + 1)^{c+d+1}$.
\end{lemma}

\begin{lemma}[Compiling a certificate language into an NDTM]
\lean{Complexity.b2_compile}
\label{Complexity.b2_compile}
\leanok
\uses{Complexity.acceptsWithin_iff_of_halts, Complexity.b2Host, Complexity.b2_host_bound, Complexity.b2_host_contract, Complexity.b2_unary_first, Complexity.e3_split_none_iff, Turing.FinNDTM, Turing.FinNDTM.AcceptsWithin, Turing.FinNDTM.DecidesInTime, Turing.FinTM, Turing.FinTM.DecidesInTime, Turing.FinTM.computesFunInTime_polyUnary, Turing.MultiTapeTM.indicator, Turing.NDTM.HaltsWithin, Turing.NDTM.HaltsWithin.mono}
\difficulty{5}
Let $L$ be a language over the binary alphabet such that $x \in L$ exactly when some binary string $u$ with $|u| = C\,(|x|+1)^c$ has $x\,u \in V$, and let $M$ be a deterministic machine deciding the verifier language $V$ within time $A(m+1)^d$. Then there are constants $K, r$ and a nondeterministic machine deciding $L$ within time $K\,(n + C(n+1)^c + 1)^r$.
\end{lemma}

\begin{theorem}[NP lies in the union of polynomial NTIME classes]
\lean{Complexity.NP_subset_iUnion_NTIME}
\label{Complexity.NP_subset_iUnion_NTIME}
\leanok
\uses{Complexity.NP, Complexity.NTIME, Complexity.b2_compile, Complexity.cont_guess_normalize, Complexity.mem_P_iff}
\difficulty{2}
Every language in $\mathrm{NP}$ (given by a polynomial-time verifier with certificates of exact polynomial length) belongs to $\bigcup_{c} \mathrm{NTIME}(n^c + 1)$.
\end{theorem}

\begin{theorem}[Machine characterization of NP]
\lean{Complexity.NP_eq_iUnion_NTIME}
\label{Complexity.NP_eq_iUnion_NTIME}
\leanok
\uses{Complexity.NP, Complexity.NP_subset_iUnion_NTIME, Complexity.NTIME, Complexity.ntime_poly_subset_NP}
\difficulty{1}
The class $\mathrm{NP}$, defined by polynomial-time verifiers with certificates of exact polynomial length, equals the union $\bigcup_{c} \mathrm{NTIME}(n^c + 1)$ of the fixed-degree nondeterministic time classes.
\end{theorem}

\subsection{The exponential split layer}

\begin{lemma}[Strict monotonicity of the exponential length map]
\lean{Complexity.e3_split_strictMono}
\label{Complexity.e3_split_strictMono}
\leanok
\difficulty{2}
For all natural numbers $C$ and $c$, the map $n \mapsto n + C \cdot 2^{(n+1)^c}$ on the natural numbers is strictly increasing.
\end{lemma}

\begin{lemma}[Uniqueness of the exponential split]
\lean{Complexity.e3_split_unique}
\label{Complexity.e3_split_unique}
\leanok
\uses{Complexity.e3Split, Complexity.e3_split_strictMono}
\difficulty{3}
Let $C, c$ be natural numbers and $x, u, y, v$ binary strings with $|u| = C\cdot 2^{(|x|+1)^c}$ and $|v| = C \cdot 2^{(|y|+1)^c}$. If $x\,u = y\,v$ as strings, then $x = y$ and $u = v$.
\end{lemma}

\begin{definition}[Bounded search for the exponential split]
\lean{Complexity.e3Split}
\label{Complexity.e3Split}
\leanok
For all natural numbers $C, c$ and $m$, the search scans all $n \le m$ and returns the least $n$ with $n + C\cdot 2^{(n+1)^c} = m$, or an explicit failure value when no such $n$ exists.
\end{definition}

\begin{lemma}[Soundness of the exponential split search]
\lean{Complexity.e3_split_spec}
\label{Complexity.e3_split_spec}
\leanok
\uses{Complexity.e3Split}
\difficulty{2}
For natural numbers $C, c, m, n$: if the exponential split search at length $m$ returns $n$, then $n \le m$ and $n + C \cdot 2^{(n+1)^c} = m$.
\end{lemma}

\begin{lemma}[Failure of the exponential split search]
\lean{Complexity.e3_split_none_iff}
\label{Complexity.e3_split_none_iff}
\leanok
\uses{Complexity.e3Split}
\difficulty{3}
For natural numbers $C, c$ and $m$, the exponential split search at length $m$ fails if and only if no natural number $n$ satisfies $n + C \cdot 2^{(n+1)^c} = m$.
\end{lemma}

\begin{definition}[Choice-word verifier with exponential width]
\lean{Complexity.e3ChoiceVerifier}
\label{Complexity.e3ChoiceVerifier}
\leanok
\uses{Turing.FinNDTM, Turing.NDTM.initCfg, Turing.NDTM.runWith}
Given a nondeterministic Turing machine $N$ over the binary alphabet and natural numbers $C, c$, the exponential choice-word verifier language consists of all strings $x\,u$ (concatenation) with $|u| = C\cdot 2^{(|x|+1)^c}$ such that the run of $N$ on $x$ under choice word $u$ has halted with output the one-bit string $1$.
\end{definition}

\begin{lemma}[Membership of a correctly split exponential word]
\lean{Complexity.e3_choice_append}
\label{Complexity.e3_choice_append}
\leanok
\uses{Complexity.e3ChoiceVerifier, Complexity.e3_split_unique, Turing.FinNDTM, Turing.NDTM.initCfg, Turing.NDTM.runWith}
\difficulty{3}
Let $N$ be a nondeterministic Turing machine over the binary alphabet, $C, c$ natural numbers, and $x, u$ binary strings with $|u| = C \cdot 2^{(|x|+1)^c}$. Then $x\,u$ lies in the exponential choice-word verifier language of $N$ with parameters $C, c$ if and only if the run of $N$ on $x$ under $u$ has halted with output the one-bit string $1$.
\end{lemma}

\begin{lemma}[Rejection without an exponential split]
\lean{Complexity.e3_choice_no_split}
\label{Complexity.e3_choice_no_split}
\leanok
\uses{Complexity.e3ChoiceVerifier, Complexity.e3Split, Complexity.e3_split_none_iff, Turing.FinNDTM}
\difficulty{3}
Let $N$ be a nondeterministic Turing machine and $C, c$ natural numbers. If the exponential split search fails at the length of a string $y$, then $y$ is not in the exponential choice-word verifier language of $N$ with parameters $C, c$.
\end{lemma}

\begin{lemma}[The padded exponential covers the budget]
\lean{Complexity.e3_choice_budget}
\label{Complexity.e3_choice_budget}
\leanok
\difficulty{1}
For all natural numbers $a, c$ and $n$, the inequality $a \cdot 2^{n^c} \le a \cdot 2^{(n+1)^c}$ holds.
\end{lemma}

\begin{lemma}[Exponential certificates for the verifier]
\lean{Complexity.e3_choice_certificate}
\label{Complexity.e3_choice_certificate}
\leanok
\uses{Complexity.acceptsWithin_iff_of_halts, Complexity.e3ChoiceVerifier, Complexity.e3_choice_append, Complexity.e3_choice_budget, Turing.FinNDTM, Turing.FinNDTM.DecidesInTime}
\difficulty{3}
Let $N$ be a nondeterministic Turing machine deciding a language $L$ within time $a \cdot 2^{n^c}$ on inputs of length $n$. For every binary string $x$: $x \in L$ if and only if there is a string $u$ with $|u| = a \cdot 2^{(|x|+1)^c}$ such that $x\,u$ lies in the exponential choice-word verifier language of $N$ with parameters $a, c$.
\end{lemma}

\begin{lemma}[Binary expansion of a shifted coefficient]
\lean{Complexity.e3_bits_shift}
\label{Complexity.e3_bits_shift}
\leanok
\difficulty{4}
For natural numbers $C \ne 0$ and $p$, the little-endian binary expansion of $C \cdot 2^p$ equals $p$ zero bits followed by the binary expansion of $C$.
\end{lemma}

\begin{lemma}[Unconditional bit-length bound for the width]
\lean{Complexity.e3_bits_length_bound}
\label{Complexity.e3_bits_length_bound}
\leanok
\uses{Complexity.e3_bits_shift}
\difficulty{3}
For all natural numbers $C, c, n, m$ with $n \le m$, the binary expansion of $C \cdot 2^{(n+1)^c}$ has length at most $(m+1)^c$ plus the length of the binary expansion of $C$.
\end{lemma}

\begin{definition}[Zero-prefixing emitter for a fixed word]
\lean{Complexity.e3ShiftTM}
\label{Complexity.e3ShiftTM}
\leanok
\uses{Turing.FinTM, Turing.FinTM.controlAction, Turing.FinTM.emitAction}
For a fixed binary word $w$, a deterministic machine with no work tapes is defined that outputs one $0$ bit per input symbol scanned and then, at the input blank, emits the fixed word $w$ and halts.
\end{definition}

\begin{lemma}[Scan invariant of the zero emitter]
\lean{Complexity.e3_shift_scan}
\label{Complexity.e3_shift_scan}
\leanok
\uses{Complexity.e3ShiftTM, Turing.Action.apply, Turing.Cfg.workTapeSymbols, Turing.MultiTapeTM.initCfg, Turing.MultiTapeTM.runFrom, Turing.MultiTapeTM.runFrom_succ_eq_step', Turing.MultiTapeTM.step, Turing.inputSymbolInner, Turing.moveInputPos_pos_of_ne_right}
\difficulty{5}
For every fixed word $w$, every input $x$ and every $t \le |x|$: after $t$ steps from the initial configuration, the zero-prefixing emitter for $w$ is still in its scanner state, its input head is at position $t + 1$, and its output so far is $t$ zero bits.
\end{lemma}

\begin{lemma}[Exact output of the zero emitter]
\lean{Complexity.e3_shift_computes}
\label{Complexity.e3_shift_computes}
\leanok
\uses{Complexity.e3ShiftTM, Complexity.e3_shift_scan, Turing.Action.apply, Turing.Cfg.inputSymbol, Turing.FinTM.ComputesInTime, Turing.FinTM.computesInTime_iff, Turing.FinTM.controlAction, Turing.FinTM.emit_halts, Turing.MultiTapeTM.initCfg, Turing.MultiTapeTM.runFrom, Turing.MultiTapeTM.runFrom_add, Turing.MultiTapeTM.runFrom_succ_eq_step', Turing.MultiTapeTM.step}
\difficulty{5}
For every fixed word $w$ and every input $x$, the zero-prefixing emitter for $w$ halts within $|x| + |w| + 2$ steps, and its complete output is the word of $|x|$ zero bits followed by $w$.
\end{lemma}

\begin{lemma}[Linear budget of the zero emitter]
\lean{Complexity.e3_shift_timed}
\label{Complexity.e3_shift_timed}
\leanok
\uses{Complexity.e3ShiftTM, Complexity.e3_shift_computes, Turing.FinTM.ComputesFunInTime, Turing.FinTM.ComputesInTime.mono}
\difficulty{2}
For every fixed word $w$, the zero-prefixing emitter computes, on each input $x$, the word of $|x|$ zero bits followed by $w$; its running time is at most the monotone linear budget $(|w|+2)(n+1)$ on inputs of length $n$.
\end{lemma}

\begin{lemma}[Polynomial-time binary evaluation of the width]
\lean{Complexity.e3_exp_bits_timed}
\label{Complexity.e3_exp_bits_timed}
\leanok
\uses{Complexity.e3_bits_shift, Complexity.e3_shift_timed, Turing.FinTM, Turing.FinTM.ComputesFunInTime, Turing.FinTM.ComputesInTime.mono, Turing.FinTM.computesFunInTime_comp, Turing.FinTM.computesFunInTime_const, Turing.FinTM.computesFunInTime_polyUnary}
\difficulty{5}
For all natural numbers $C$ and $c$ there exist a deterministic machine and a constant $A$ computing, on every input $x$, the little-endian binary expansion of $C \cdot 2^{(|x|+1)^c}$ within $A\,(|x|+1)^{c+1}$ steps.
\end{lemma}

\begin{definition}[Split emitter for the exponential length equation]
\lean{Complexity.e3SplitWord}
\label{Complexity.e3SplitWord}
\leanok
\uses{Complexity.e3Split, Turing.pairEncode}
For natural numbers $C, c$ and a binary string $y$, the emitted word is the self-delimiting encoding of the pair formed by the first $i$ bits of $y$ and the remaining bits of $y$, when the exponential split search returns $i$ at length $|y|$; it is the empty word on failure.
\end{definition}

\begin{lemma}[Linear bound on the emitted exponential split]
\lean{Complexity.e3_split_length}
\label{Complexity.e3_split_length}
\leanok
\uses{Complexity.e3Split, Complexity.e3SplitWord, Complexity.e3_split_spec, Turing.pairEncode}
\difficulty{3}
For all natural numbers $C, c$ and every binary string $y$, the word emitted by the exponential split emitter on $y$ has length at most $2|y| + 2$.
\end{lemma}

\begin{lemma}[The loader decides the exponential verifier]
\lean{Complexity.e3_split_answer}
\label{Complexity.e3_split_answer}
\leanok
\uses{Complexity.contPairTM, Complexity.cont_pair_computes, Complexity.cont_pair_empty, Complexity.e3ChoiceVerifier, Complexity.e3Split, Complexity.e3SplitWord, Complexity.e3_choice_append, Complexity.e3_choice_no_split, Complexity.e3_split_length, Complexity.e3_split_spec, Turing.FinNDTM, Turing.FinTM.ComputesInTime, Turing.FinTM.ComputesInTime.mono, Turing.MultiTapeTM.indicator, Turing.NDTM.initCfg, Turing.NDTM.runWith}
\difficulty{5}
Let $N$ be a nondeterministic Turing machine and $C, c$ natural numbers. For every binary string $y$, the pair-loader machine of $N$, run on the word emitted by the exponential split emitter for $y$, halts within $3(2|y|+3)$ steps, and its output is the indicator bit of membership of $y$ in the exponential choice-word verifier language of $N$.
\end{lemma}

\begin{lemma}[Polynomial split recovery puts the verifier in P]
\lean{Complexity.e3_verifier_of_split}
\label{Complexity.e3_verifier_of_split}
\leanok
\uses{Complexity.P, Complexity.contPairTM, Complexity.e3ChoiceVerifier, Complexity.e3SplitStep, Complexity.e3SplitWord, Complexity.e3_split_answer, Complexity.e3_split_length, Complexity.mem_P_iff, Turing.Cfg.init, Turing.FinNDTM, Turing.FinTM, Turing.FinTM.ComputesFunInTime, Turing.FinTM.ComputesInTime, Turing.FinTM.ComputesInTime.mono, Turing.FinTM.VirtualTag, Turing.FinTM.bufferedCompTM, Turing.FinTM.bufferedComp_start, Turing.FinTM.bufferedSecondCfg, Turing.FinTM.bufferedSecondCfg_run, Turing.FinTM.computesInTime_iff, Turing.MultiTapeTM.indicator, Turing.MultiTapeTM.initCfg, Turing.MultiTapeTM.runFrom_add}
\difficulty{5}
Let $N$ be a nondeterministic Turing machine and $C, c$ natural numbers. If some deterministic machine computes the exponential split-emitter function within time $A(n+1)^{r+1}$ for constants $A, r$, then the exponential choice-word verifier language of $N$ with parameters $C, c$ belongs to $\mathrm{P}$.
\end{lemma}

\begin{definition}[One round of the bounded unary search]
\lean{Complexity.e3SplitStep}
\label{Complexity.e3SplitStep}
\leanok
Given an input word $w$ and a search state $s$ (a unary candidate), one round appends a $1$ bit to $s$ when $|s| \le |w|$ and leaves $s$ unchanged once it is past the end of the input.
\end{definition}

\begin{definition}[Acceptance test of a search round]
\lean{Complexity.e3SplitAccept}
\label{Complexity.e3SplitAccept}
\leanok
For all natural numbers $C, c$, the round test on an input word $w$ and a search state $s$ returns true exactly when $|s| + C \cdot 2^{(|s|+1)^c} = |w|$.
\end{definition}

\begin{lemma}[Invariance of the search-state bound]
\lean{Complexity.e3_step_inv}
\label{Complexity.e3_step_inv}
\leanok
\uses{Complexity.e3SplitStep}
\difficulty{2}
For all words $w$ and $s$ with $|s| \le |w| + 1$, the state produced by one search round still has length at most $|w| + 1$.
\end{lemma}

\begin{lemma}[The search orbit is unary]
\lean{Complexity.e3_step_orbit}
\label{Complexity.e3_step_orbit}
\leanok
\uses{Complexity.e3SplitStep}
\difficulty{3}
For every word $w$ and every $i \le |w| + 1$, iterating the search round $i$ times from the empty state yields the all-ones word of length $i$.
\end{lemma}

\begin{lemma}[First-match search respects pointwise equality]
\lean{Complexity.e3_find_congr}
\label{Complexity.e3_find_congr}
\leanok
\difficulty{3}
Let $p$ and $q$ be Boolean predicates that agree on every element of a given list. Then the first element of the list satisfying $p$ equals the first element satisfying $q$.
\end{lemma}

\begin{lemma}[The orbit search equals the split search]
\lean{Complexity.e3_find_eq}
\label{Complexity.e3_find_eq}
\leanok
\uses{Complexity.e3Split, Complexity.e3SplitAccept, Complexity.e3SplitStep, Complexity.e3_find_congr, Complexity.e3_step_orbit}
\difficulty{3}
For all natural numbers $C, c$ and every word $w$, searching the indices $i \le |w|$ for the first round state (the $i$-fold iterate from the empty state) that passes the round test returns exactly the result of the exponential split search at length $|w|$.
\end{lemma}

\begin{lemma}[Loop payload agrees with the split emitter]
\lean{Complexity.e3_loop_result}
\label{Complexity.e3_loop_result}
\leanok
\uses{Complexity.e3Split, Complexity.e3SplitAccept, Complexity.e3SplitStep, Complexity.e3SplitWord, Complexity.e3_find_eq, Complexity.e3_split_spec, Complexity.e3_step_orbit, Turing.pairEncode}
\difficulty{3}
For natural numbers $C, c$ and a word $w$: taking the first accepted index of the bounded orbit search and emitting the encoded pair of the corresponding prefix and suffix of $w$ (or the empty word when the search is exhausted) yields exactly the word of the exponential split emitter on $w$.
\end{lemma}

\begin{lemma}[Budget growth of the result-bearing loop]
\lean{Complexity.e3_loop_bound}
\label{Complexity.e3_loop_bound}
\leanok
\difficulty{3}
For all natural numbers $b, A, r$ and $n$, the product $b\,(A(n+1)^{r+1} + 1)(n+2)$ is at most $2b(A+1)\,(n+1)^{r+2}$.
\end{lemma}

\begin{lemma}[A loop body yields polynomial split recovery]
\lean{Complexity.e3_split_of_body}
\label{Complexity.e3_split_of_body}
\leanok
\uses{Complexity.e3SplitAccept, Complexity.e3SplitStep, Complexity.e3SplitWord, Complexity.e3_loop_bound, Complexity.e3_loop_result, Complexity.e3_step_inv, Turing.Cfg.ofWords, Turing.FinTM, Turing.FinTM.ComputesFunInTime, Turing.FinTM.ComputesInTime.mono, Turing.FinTM.computesFunInTime_lengthBits, Turing.FinTM.exists_loopFindTM, Turing.MultiTapeTM.initCfg, Turing.MultiTapeTM.runFrom, Turing.pairEncode, Turing.stateWord}
\difficulty{5}
Fix natural numbers $C, c$ and suppose a deterministic loop-body machine with a designated anchor state satisfies two timed contracts within $A(|w|+1)^{r+1}$ steps on every input $w$: a startup contract that first reaches the anchor with clean scratch, and a round contract that, from any anchored search state, either halts with the encoded pair payload when the round test accepts, or returns to the anchor with the stepped search state. Then some deterministic machine computes the exponential split-emitter function within $B(n+1)^{r+2}$ steps for a constant $B$.
\end{lemma}

\begin{lemma}[Prefix equality extended by one position]
\lean{Complexity.e3c_take_succ_eq}
\label{Complexity.e3c_take_succ_eq}
\leanok
\difficulty{3}
For all binary strings $u, v$ and every index $j$, the length-$(j+1)$ prefixes of $u$ and $v$ agree if and only if their length-$j$ prefixes agree and their optional bits at index $j$ agree, a missing bit counting as different from any present bit.
\end{lemma}

\begin{lemma}[Injectivity of the binary expansion]
\lean{Complexity.e3c_bits_injective}
\label{Complexity.e3c_bits_injective}
\leanok
\difficulty{4}
The map sending a natural number to its little-endian binary expansion is injective: two numbers with the same bit list are equal.
\end{lemma}

\begin{theorem}[Fixed-exponent nondeterministic time lies in NEXP]
\lean{Complexity.ntime_expPow_subset_NEXP}
\label{Complexity.ntime_expPow_subset_NEXP}
\leanok
\uses{Complexity.NEXP, Complexity.NTIME, Complexity.e3ChoiceVerifier, Complexity.e3Split, Complexity.e3SplitWord, Complexity.e3_choice_certificate, Complexity.e3_exp_bits_timed, Complexity.e3_verifier_of_split, Turing.FinTM.ComputesInTime, Turing.FinTM.ComputesInTime.mono, Turing.FinTM.computesFunInTime_splitSolveWith, Turing.solveSplitWith}
\difficulty{5}
For every natural number $c$, the class $\mathrm{NTIME}(2^{n^c})$ is contained in the certificate-form class $\mathrm{NEXP}$: every language decided by some nondeterministic Turing machine within $a \cdot 2^{n^c}$ steps admits a polynomial-time verifier with certificates of exact length $a \cdot 2^{(n+1)^c}$.
\end{theorem}

\subsection{The binary countdown and the exponential compilation}

\begin{definition}[Fixed-width little-endian decrement]
\lean{Complexity.a2Debit}
\label{Complexity.a2Debit}
\leanok
For a binary word viewed as a little-endian counter, the decrement returns a word of the same width together with a success flag: a leading $1$ is cleared (after turning the traversed low zero bits into ones), and on underflow, when the word has value zero, the flag is false and the width is still unchanged.
\end{definition}

\begin{definition}[Length of the borrow scan]
\lean{Complexity.a2BorrowPos}
\label{Complexity.a2BorrowPos}
\leanok
For a binary word viewed as a little-endian counter, the borrow position is the number of leading $0$ bits, that is, the cells a decrement must traverse before finding a $1$ or the end of the word.
\end{definition}

\begin{lemma}[The borrow scan stays inside the word]
\lean{Complexity.a2BorrowPos_le}
\label{Complexity.a2BorrowPos_le}
\leanok
\uses{Complexity.a2BorrowPos}
\difficulty{3}
For every binary word $u$, the borrow position of $u$ is at most $|u|$.
\end{lemma}

\begin{lemma}[Width preservation of the decrement]
\lean{Complexity.a2Debit_length}
\label{Complexity.a2Debit_length}
\leanok
\uses{Complexity.a2Debit}
\difficulty{3}
For every binary word $u$, the word part of the fixed-width decrement of $u$ has the same length as $u$, whether or not the decrement succeeds.
\end{lemma}

\begin{definition}[Little-endian value of a counter word]
\lean{Complexity.a2Value}
\label{Complexity.a2Value}
\leanok
The value of a binary word read as a little-endian counter is defined recursively: the empty word has value $0$, and a word with first bit $b$ and tail $t$ has value twice the value of $t$, plus one exactly when $b$ is $1$.
\end{definition}

\begin{lemma}[The binary expansion has its declared value]
\lean{Complexity.a2Value_bits}
\label{Complexity.a2Value_bits}
\leanok
\uses{Complexity.a2Value}
\difficulty{4}
For every natural number $n$, the little-endian value of the binary expansion of $n$ equals $n$.
\end{lemma}

\begin{lemma}[Arithmetic correctness of the decrement]
\lean{Complexity.a2Debit_value}
\label{Complexity.a2Debit_value}
\leanok
\uses{Complexity.a2Debit, Complexity.a2Value}
\difficulty{4}
For every binary word $u$: if the fixed-width decrement of $u$ succeeds, the value of the resulting word plus one equals the value of $u$; if it fails, the value of $u$ is zero.
\end{lemma}

\begin{lemma}[Buffer read at a suffix boundary]
\lean{Complexity.a2Buffer_read}
\label{Complexity.a2Buffer_read}
\leanok
\uses{Complexity.a2DebitTM, Turing.FinTM.bufferTape, Turing.FinTM.bufferTape_nat}
\difficulty{2}
For all binary words $p$ and $s$, reading the buffer tape holding $p\,s$ at cell $|p|$ returns the first bit of $s$, or blank when $s$ is empty.
\end{lemma}

\begin{lemma}[Buffer write at a suffix boundary]
\lean{Complexity.a2Buffer_write}
\label{Complexity.a2Buffer_write}
\leanok
\uses{Complexity.a2DebitTM, Turing.FinTM.bufferTape}
\difficulty{4}
For all binary words $p$ and $s$ and all bits $b, b'$, overwriting cell $|p|$ of the buffer tape holding $p\,b\,s$ with the bit $b'$ yields exactly the buffer tape holding $p\,b'\,s$.
\end{lemma}

\begin{definition}[Native binary countdown emitter]
\lean{Complexity.a2DebitTM}
\label{Complexity.a2DebitTM}
\leanok
\uses{Turing.FinTM, Turing.FinTM.controlAction}
A deterministic machine with one work tape is defined that copies its binary input onto the work tape, rewinds, and then repeatedly performs a fixed-width little-endian decrement with borrow and rewind phases, emitting one $1$ bit after each successful decrement; underflow halts silently.
\end{definition}

\begin{definition}[Borrow-phase configuration of the countdown]
\lean{Complexity.a2DebitCfg}
\label{Complexity.a2DebitCfg}
\leanok
\uses{Complexity.a2DebitTM, Turing.Cfg, Turing.FinTM.bufferTape}
Given an already-emitted output word, a physical input with an arbitrary head position, a borrow-phase control state, a work-head value and a counter word, the corresponding configuration of the countdown emitter holds the counter on its single work tape and the emitted word as physical output.
\end{definition}

\begin{lemma}[One borrow transition]
\lean{Complexity.a2Borrow_step}
\label{Complexity.a2Borrow_step}
\leanok
\uses{Complexity.a2Buffer_read, Complexity.a2Buffer_write, Complexity.a2DebitCfg, Complexity.a2DebitTM, Turing.Action.apply, Turing.Cfg.workTapeSymbols, Turing.MultiTapeTM.step, Turing.moveInputPos_zero}
\difficulty{5}
For every emitted prefix, every input and position, and every counter word split as $p\,s$: one step of the countdown emitter in the borrow state with work head at cell $|p|$ clears a leading $1$ of $s$ and turns to rewind with success, converts a leading $0$ of $s$ to $1$ and advances, or, at the blank end of the word, turns to rewind with failure; in each case the width is preserved and nothing is emitted.
\end{lemma}

\begin{lemma}[Exact cost of the borrow phase]
\lean{Complexity.a2Borrow_run}
\label{Complexity.a2Borrow_run}
\leanok
\uses{Complexity.a2BorrowPos, Complexity.a2Borrow_step, Complexity.a2Debit, Complexity.a2DebitCfg, Complexity.a2DebitTM, Turing.MultiTapeTM.runFrom, Turing.MultiTapeTM.runFrom_succ_eq_step}
\difficulty{5}
For every counter word $u$ and every prefix $p$: starting in the borrow state with work head at cell $|p|$ over the word $p\,u$, exactly $\beta + 1$ steps of the countdown emitter, where $\beta$ is the borrow position of $u$, reach the rewind state carrying the success flag of the decrement of $u$, with the work tape now holding $p$ followed by the decremented word.
\end{lemma}

\begin{lemma}[Exact cost of the counter rewind]
\lean{Complexity.a2Borrow_rewind}
\label{Complexity.a2Borrow_rewind}
\leanok
\uses{Complexity.a2DebitCfg, Complexity.a2DebitTM, Turing.Action.apply, Turing.Cfg.workTapeSymbols, Turing.FinTM.bufferTape_left, Turing.FinTM.bufferTape_nat, Turing.MultiTapeTM.runFrom, Turing.MultiTapeTM.runFrom_succ_eq_step, Turing.MultiTapeTM.runFrom_zero, Turing.MultiTapeTM.step, Turing.moveInputPos_zero}
\difficulty{4}
For every counter word $u$, every flag $b$ and every $j \le |u|$: from the rewind state with work head at cell $j-1$, exactly $j+1$ steps of the countdown emitter reach the round-return state carrying $b$, with head at cell zero and the counter unchanged.
\end{lemma}

\begin{lemma}[Complete decrement round]
\lean{Complexity.a2Borrow_correct}
\label{Complexity.a2Borrow_correct}
\leanok
\uses{Complexity.a2BorrowPos, Complexity.a2BorrowPos_le, Complexity.a2Borrow_rewind, Complexity.a2Borrow_run, Complexity.a2Debit, Complexity.a2DebitCfg, Complexity.a2DebitTM, Complexity.a2Debit_length, Turing.MultiTapeTM.runFrom, Turing.MultiTapeTM.runFrom_add}
\difficulty{4}
For every counter word $u$: a full borrow-and-rewind round of the countdown emitter from cell zero costs $2\beta + 2 \le 2|u| + 2$ steps, where $\beta$ is the borrow position of $u$, and ends at cell zero in the round-return state carrying the success flag, with the counter replaced by its fixed-width decrement and no output emitted.
\end{lemma}

\begin{lemma}[The countdown emits exactly its value]
\lean{Complexity.a2_countdown}
\label{Complexity.a2_countdown}
\leanok
\uses{Complexity.a2BorrowPos, Complexity.a2Borrow_correct, Complexity.a2Debit, Complexity.a2DebitCfg, Complexity.a2DebitTM, Complexity.a2Debit_length, Complexity.a2Debit_value, Complexity.a2Value, Turing.Action.apply, Turing.FinTM.controlAction, Turing.MultiTapeTM.runFrom, Turing.MultiTapeTM.runFrom_add, Turing.MultiTapeTM.step, Turing.moveInputPos_zero}
\difficulty{6}
For every input, head position and counter word $u$ of value $n$, and every already-emitted word: within $(n+1)(2|u|+3)$ steps from the borrow state at cell zero, the countdown emitter halts with output equal to the already-emitted word followed by exactly $n$ ones.
\end{lemma}

\begin{definition}[Startup configuration of the countdown]
\lean{Complexity.a2LoadCfg}
\label{Complexity.a2LoadCfg}
\leanok
\uses{Complexity.a2DebitTM, Turing.Cfg, Turing.FinTM.bufferTape}
Given a copy-phase flag, a physical input position, a copied word and a work-head value, the corresponding configuration of the countdown emitter is in its startup control state with the copied word on the work tape and empty output.
\end{definition}

\begin{lemma}[One input-copy step of the countdown]
\lean{Complexity.a2_copy_step}
\label{Complexity.a2_copy_step}
\leanok
\uses{Complexity.a2DebitTM, Complexity.a2LoadCfg, Turing.Action.apply, Turing.Cfg.inputSymbol, Turing.FinTM.bufferTape_append, Turing.MultiTapeTM.step, Turing.inputSymbolInner, Turing.moveInputPos, Turing.moveInputPos_pos_of_ne_right}
\difficulty{4}
For every input $x$ and every index $i < |x|$, one step of the countdown emitter from the copy configuration holding the first $i$ bits of $x$ appends the bit at index $i$ to the work tape and advances both heads, emitting nothing.
\end{lemma}

\begin{lemma}[Startup copies the whole input]
\lean{Complexity.a2_copy_run}
\label{Complexity.a2_copy_run}
\leanok
\uses{Complexity.a2DebitTM, Complexity.a2LoadCfg, Complexity.a2_copy_step, Turing.Cfg.init, Turing.FinTM.bufferTape, Turing.MultiTapeTM.initCfg, Turing.MultiTapeTM.runFrom, Turing.MultiTapeTM.runFrom_succ_eq_step'}
\difficulty{3}
For every input $x$ and every $i \le |x|$, the $i$-step run of the countdown emitter from its blank-tape initial configuration is the copy configuration holding the first $i$ bits of $x$ with input head at position $i+1$.
\end{lemma}

\begin{lemma}[Startup rewind of the copied counter]
\lean{Complexity.a2_load_rewind}
\label{Complexity.a2_load_rewind}
\leanok
\uses{Complexity.a2DebitCfg, Complexity.a2DebitTM, Complexity.a2LoadCfg, Turing.Action.apply, Turing.Cfg.workTapeSymbols, Turing.FinTM.bufferTape_left, Turing.FinTM.bufferTape_nat, Turing.MultiTapeTM.runFrom, Turing.MultiTapeTM.runFrom_succ_eq_step, Turing.MultiTapeTM.runFrom_zero, Turing.MultiTapeTM.step, Turing.moveInputPos_zero}
\difficulty{4}
For every input position, every counter word $u$ and every $j \le |u|$: from the startup-rewind configuration with work head at cell $j-1$, exactly $j+1$ steps of the countdown emitter reach the borrow state at cell zero with counter $u$ and empty output.
\end{lemma}

\begin{lemma}[Complete startup of the countdown]
\lean{Complexity.a2_start}
\label{Complexity.a2_start}
\leanok
\uses{Complexity.a2DebitCfg, Complexity.a2DebitTM, Complexity.a2LoadCfg, Complexity.a2_copy_run, Complexity.a2_load_rewind, Turing.Action.apply, Turing.Cfg.inputSymbol, Turing.MultiTapeTM.initCfg, Turing.MultiTapeTM.runFrom, Turing.MultiTapeTM.runFrom_add, Turing.MultiTapeTM.runFrom_succ_eq_step', Turing.MultiTapeTM.step, Turing.moveInputPos_zero}
\difficulty{4}
For every input $x$, the $(2|x| + 2)$-step run of the countdown emitter from its blank-tape initial configuration is the borrow-state configuration at cell zero with the full input $x$ as counter, input head at the right boundary and empty output.
\end{lemma}

\begin{lemma}[Timed decoding contract of the countdown]
\lean{Complexity.a2_decode_computes}
\label{Complexity.a2_decode_computes}
\leanok
\uses{Complexity.a2DebitTM, Complexity.a2Value, Complexity.a2_countdown, Complexity.a2_start, Turing.FinTM.ComputesInTime, Turing.FinTM.ComputesInTime.mono, Turing.FinTM.computesInTime_iff, Turing.MultiTapeTM.runFrom_add}
\difficulty{3}
For every binary input $x$, the countdown emitter computes the all-ones word whose length is the little-endian value $v$ of $x$, within $2|x| + 2 + (v+1)(2|x|+3)$ steps.
\end{lemma}

\begin{lemma}[Exponential unary scheduler via binary countdown]
\lean{Complexity.a2_exp_scheduler}
\label{Complexity.a2_exp_scheduler}
\leanok
\uses{Complexity.a2DebitTM, Complexity.a2Value_bits, Complexity.a2_decode_computes, Complexity.e3_exp_bits_timed, Turing.Cfg.init, Turing.FinTM, Turing.FinTM.ComputesFunInTime, Turing.FinTM.ComputesInTime, Turing.FinTM.ComputesInTime.mono, Turing.FinTM.VirtualTag, Turing.FinTM.bufferedCompTM, Turing.FinTM.bufferedComp_start, Turing.FinTM.bufferedSecondCfg, Turing.FinTM.bufferedSecondCfg_run, Turing.FinTM.computesInTime_iff, Turing.MultiTapeTM.initCfg, Turing.MultiTapeTM.output_length_le, Turing.MultiTapeTM.runFrom_add}
\difficulty{6}
For all natural numbers $C$ and $c$ there exist a deterministic machine $S$ and a constant $B$ such that on every input $x$, $S$ outputs the all-ones word of length $C \cdot 2^{(|x|+1)^c}$ within $B\,(|x| + C\cdot 2^{(|x|+1)^c} + 1)^{c+2}$ steps.
\end{lemma}

\begin{lemma}[Envelope of the exponential host ledger]
\lean{Complexity.a2_host_bound}
\label{Complexity.a2_host_bound}
\leanok
\difficulty{4}
For all natural numbers $Q, B, r, A, d, n$ and $T$ with $T \le B(n+Q+1)^r$: the total $2n + 2 + T + (n + Q + 2 + (A(n+Q+1)^d + 1))$ is at most $(B+A+5)(n+Q+1)^{r+d+1}$.
\end{lemma}

\begin{lemma}[Compiling an exponential certificate language]
\lean{Complexity.a2_compile}
\label{Complexity.a2_compile}
\leanok
\uses{Complexity.a2_exp_scheduler, Complexity.a2_host_bound, Complexity.acceptsWithin_iff_of_halts, Complexity.b2Host, Complexity.b2_host_contract, Complexity.b2_unary_first, Turing.FinNDTM, Turing.FinNDTM.AcceptsWithin, Turing.FinNDTM.DecidesInTime, Turing.FinTM, Turing.FinTM.DecidesInTime, Turing.MultiTapeTM.indicator, Turing.NDTM.HaltsWithin, Turing.NDTM.HaltsWithin.mono}
\difficulty{5}
Let $L$ be a language over the binary alphabet such that $x \in L$ exactly when some binary string $u$ with $|u| = C \cdot 2^{(|x|+1)^c}$ has $x\,u \in V$, and let $M$ be a deterministic machine deciding the verifier language $V$ within $A(m+1)^d$. Then there are constants $K, r$ and a nondeterministic machine deciding $L$ within $K\,(n + C\cdot 2^{(n+1)^c} + 1)^r$.
\end{lemma}

\begin{lemma}[Absorbing a polynomial exponent]
\lean{Complexity.a2_exponent_bound}
\label{Complexity.a2_exponent_bound}
\leanok
\difficulty{4}
For all natural numbers $K$ and $k$ there exist constants $A$ and $e$ such that $2^{K(n+1)^k} \le A \cdot 2^{n^e}$ for every natural number $n$.
\end{lemma}

\begin{lemma}[Exponential-time envelope of the compiled budget]
\lean{Complexity.a2_envelope}
\label{Complexity.a2_envelope}
\leanok
\uses{Complexity.a2_exponent_bound}
\difficulty{4}
For all natural numbers $C, c, K, r$ there exist constants $A$ and $e$ with $K\,(n + C\cdot 2^{(n+1)^c} + 1)^r \le A \cdot 2^{n^e}$ for every natural number $n$.
\end{lemma}

\begin{lemma}[Landing a compiled exponential decider in NTIME]
\lean{Complexity.a2_normalize}
\label{Complexity.a2_normalize}
\leanok
\uses{Complexity.NTIME, Complexity.a2_envelope, Complexity.acceptsWithin_iff_of_halts, Turing.FinNDTM, Turing.FinNDTM.DecidesInTime, Turing.NDTM.HaltsWithin.mono}
\difficulty{3}
Let $L$ be a language over the binary alphabet and $C, c$ natural numbers. If some nondeterministic machine decides $L$ within time $K(n + C\cdot 2^{(n+1)^c} + 1)^r$ for constants $K, r$, then $L$ belongs to $\bigcup_e \mathrm{NTIME}(2^{n^e})$.
\end{lemma}

\begin{theorem}[NEXP lies in the union of exponential NTIME classes]
\lean{Complexity.NEXP_subset_iUnion_NTIME}
\label{Complexity.NEXP_subset_iUnion_NTIME}
\leanok
\uses{Complexity.NEXP, Complexity.NTIME, Complexity.a2_compile, Complexity.a2_normalize, Complexity.mem_P_iff}
\difficulty{2}
Every language in the certificate-form class $\mathrm{NEXP}$ (given by a polynomial-time verifier with certificates of exact exponential length) belongs to $\bigcup_c \mathrm{NTIME}(2^{n^c})$.
\end{theorem}

\begin{theorem}[The NTIME form of NEXP]
\lean{Complexity.NEXP_eq_iUnion_NTIME}
\label{Complexity.NEXP_eq_iUnion_NTIME}
\leanok
\uses{Complexity.NEXP, Complexity.NEXP_subset_iUnion_NTIME, Complexity.NTIME, Complexity.ntime_expPow_subset_NEXP}
\difficulty{1}
The certificate-form class $\mathrm{NEXP}$ equals the union $\bigcup_c \mathrm{NTIME}(2^{n^c})$ of the fixed-exponent nondeterministic time classes.
\end{theorem}

\subsection{Polynomial-time calculus, padding, and Theorem 2.22}

\begin{lemma}[Linear budgets are polynomial]
\lean{Complexity.a3n_poly_linear}
\label{Complexity.a3n_poly_linear}
\leanok
\uses{Complexity.PolyTimeComputable, Turing.FinTM, Turing.FinTM.ComputesFunInTime}
\difficulty{2}
Let $f$ be a function on binary strings. If some deterministic machine computes $f$ within a budget $a(n+1)$ on inputs of length $n$, for a constant $a$, then $f$ is polynomial-time computable.
\end{lemma}

\begin{lemma}[Constant functions are polynomial time]
\lean{Complexity.a3n_poly_const}
\label{Complexity.a3n_poly_const}
\leanok
\uses{Complexity.PolyTimeComputable, Complexity.a3n_poly_linear, Turing.FinTM.computesFunInTime_const}
\difficulty{1}
For every fixed binary word $w$, the constant function sending every input to $w$ is polynomial-time computable.
\end{lemma}

\begin{lemma}[Polynomial-time conditional branching]
\lean{Complexity.a3n_poly_cond}
\label{Complexity.a3n_poly_cond}
\leanok
\uses{Complexity.PolyTimeComputable, Turing.FinTM.ComputesInTime.mono, Turing.FinTM.computesFunInTime_cond}
\difficulty{4}
Let $p$ be a Boolean test whose single-bit indicator is polynomial-time computable, and let $f$ and $g$ be polynomial-time computable functions on binary strings. Then the function that returns $f(x)$ when $p(x)$ holds and $g(x)$ otherwise is polynomial-time computable.
\end{lemma}

\begin{definition}[Total first projection of an encoded pair]
\lean{Complexity.a3n_fst}
\label{Complexity.a3n_fst}
\leanok
\uses{Turing.pairDecode}
For a binary string $z$, the projection returns the first component of the decoded pair when $z$ is a valid self-delimiting pair encoding, and the empty word otherwise.
\end{definition}

\begin{definition}[Total second projection of an encoded pair]
\lean{Complexity.a3n_snd}
\label{Complexity.a3n_snd}
\leanok
\uses{Turing.pairDecode}
For a binary string $z$, the projection returns the second component of the decoded pair when $z$ is a valid self-delimiting pair encoding, and the empty word otherwise.
\end{definition}

\begin{definition}[Guarded concatenation of an encoded pair]
\lean{Complexity.a3n_concat}
\label{Complexity.a3n_concat}
\leanok
\uses{Turing.pairDecode}
For a binary string $z$, the result is the concatenation of the two components when $z$ is a valid self-delimiting pair encoding, and the empty word otherwise.
\end{definition}

\begin{definition}[Payload map on encoded pairs]
\lean{Complexity.a3n_map}
\label{Complexity.a3n_map}
\leanok
\uses{Turing.pairDecode, Turing.pairEncode}
Given a function $g$ on binary strings and a string $z$: when $z$ decodes as a pair $(a, b)$, the result is the encoding of the pair $(a, g(b))$, and otherwise the empty word.
\end{definition}

\begin{lemma}[The payload map preserves polynomial time]
\lean{Complexity.a3n_poly_map}
\label{Complexity.a3n_poly_map}
\leanok
\uses{Complexity.PolyTimeComputable, Complexity.a3n_map, Turing.FinTM.ComputesInTime.mono, Turing.FinTM.computesFunInTime_pairMapSnd}
\difficulty{4}
Let $g$ be a polynomial-time computable function on binary strings. Then the function that decodes its input as a pair $(a,b)$ and re-encodes the pair $(a, g(b))$, returning the empty word on malformed input, is polynomial-time computable.
\end{lemma}

\begin{lemma}[Polynomial-time pairing of two functions]
\lean{Complexity.a3n_poly_pair}
\label{Complexity.a3n_poly_pair}
\leanok
\uses{Complexity.PolyTimeComputable, Complexity.PolyTimeComputable.comp, Complexity.a3n_concat, Complexity.a3n_fst, Complexity.a3n_map, Complexity.a3n_poly_const, Complexity.a3n_poly_linear, Complexity.a3n_poly_map, Complexity.a3n_snd, Turing.FinTM.computesFunInTime_pairConcat, Turing.FinTM.computesFunInTime_pairDup, Turing.FinTM.computesFunInTime_pairFst, Turing.FinTM.computesFunInTime_pairSnd, Turing.pairDecode_pairEncode, Turing.pairEncode}
\difficulty{5}
Let $f$ and $g$ be polynomial-time computable functions on binary strings. Then the function sending $x$ to the self-delimiting encoding of the pair $(f(x), g(x))$ is polynomial-time computable.
\end{lemma}

\begin{lemma}[Timed buffered composition at the actual word]
\lean{Complexity.a3n_comp_at}
\label{Complexity.a3n_comp_at}
\leanok
\uses{Turing.Cfg.init, Turing.FinTM, Turing.FinTM.ComputesInTime, Turing.FinTM.ComputesInTime.mono, Turing.FinTM.VirtualTag, Turing.FinTM.bufferedCompTM, Turing.FinTM.bufferedComp_start, Turing.FinTM.bufferedSecondCfg, Turing.FinTM.bufferedSecondCfg_run, Turing.FinTM.computesInTime_iff, Turing.MultiTapeTM.initCfg, Turing.MultiTapeTM.runFrom_add}
\difficulty{3}
Let $F$ and $G$ be deterministic machines, and suppose $F$ computes $y$ from $x$ within $s$ steps and $G$ computes $z$ from $y$ within $t$ steps. Then the buffered composition of $F$ and $G$ computes $z$ from $x$ within $s + |y| + 2 + t$ steps.
\end{lemma}

\begin{definition}[Pair comparator machine]
\lean{Complexity.a3nEqTM}
\label{Complexity.a3nEqTM}
\leanok
\uses{Turing.FinTM}
A deterministic machine with one work tape is defined that, on a self-delimiting encoded pair, decodes the doubled-bit prefix onto its work tape, rewinds, then compares the stored word with the suffix bit by bit, including the final blank, and halts emitting a single bit recording whole-word equality.
\end{definition}

\begin{definition}[Configuration of the pair comparator]
\lean{Complexity.a3nEqCfg}
\label{Complexity.a3nEqCfg}
\leanok
\uses{Complexity.a3nEqTM, Turing.Cfg, Turing.FinTM.bufferTape}
Given a control index, a comparison flag, a physical input position, a decoded word and a work-head value, the corresponding comparator configuration holds the decoded word on its single work tape, with empty physical output.
\end{definition}

\begin{lemma}[Two comparator steps decode a doubled bit]
\lean{Complexity.a3n_eq_double}
\label{Complexity.a3n_eq_double}
\leanok
\uses{Complexity.a3nEqCfg, Complexity.a3nEqTM, Turing.Action.apply, Turing.Cfg.inputSymbol, Turing.FinTM.bufferTape_append, Turing.MultiTapeTM.runFrom, Turing.MultiTapeTM.runFrom_succ_eq_step, Turing.MultiTapeTM.runFrom_zero, Turing.MultiTapeTM.step, Turing.inputSymbolInner, Turing.moveInputPos, Turing.moveInputPos_pos_of_ne_right}
\difficulty{5}
Suppose the comparator input has the form $y = p\,b\,b\,r$ for a bit $b$ and words $p, r$. From the parser state at position $|p|+1$ with decoded word $u$, exactly two steps of the pair comparator append $b$ to the work tape and move to position $|p|+3$, emitting nothing.
\end{lemma}

\begin{lemma}[The comparator recognizes the separator]
\lean{Complexity.a3n_eq_separator}
\label{Complexity.a3n_eq_separator}
\leanok
\uses{Complexity.a3nEqCfg, Complexity.a3nEqTM, Turing.Action.apply, Turing.Cfg.inputSymbol, Turing.MultiTapeTM.runFrom, Turing.MultiTapeTM.runFrom_succ_eq_step, Turing.MultiTapeTM.runFrom_zero, Turing.MultiTapeTM.step, Turing.inputSymbolInner, Turing.moveInputPos, Turing.moveInputPos_pos_of_ne_right}
\difficulty{4}
Suppose the comparator input has the form $y = p\,0\,1\,v$ for words $p, v$. From the parser state at position $|p|+1$ with decoded word $u$, exactly two steps of the pair comparator enter the rewind state at position $|p|+3$, with the work tape and output untouched and the work head moved one cell left.
\end{lemma}

\begin{lemma}[Exact parse by the comparator]
\lean{Complexity.a3n_eq_parse}
\label{Complexity.a3n_eq_parse}
\leanok
\uses{Complexity.a3nEqCfg, Complexity.a3nEqTM, Complexity.a3n_eq_double, Complexity.a3n_eq_separator, Turing.MultiTapeTM.runFrom, Turing.MultiTapeTM.runFrom_add, Turing.pairEncode}
\difficulty{5}
Suppose the comparator input has the form $y = p\,e$ where $e$ is the self-delimiting encoding of a pair $(u, v)$. From the parser state at position $|p|+1$ with decoded word $a$, exactly $2|u|+2$ steps of the pair comparator reach the rewind state at position $|p|+2|u|+3$ with decoded word $a\,u$.
\end{lemma}

\begin{lemma}[Comparator rewind of the decoded word]
\lean{Complexity.a3n_eq_rewind}
\label{Complexity.a3n_eq_rewind}
\leanok
\uses{Complexity.a3nEqCfg, Complexity.a3nEqTM, Turing.Action.apply, Turing.Cfg.workTapeSymbols, Turing.FinTM.bufferTape_left, Turing.FinTM.bufferTape_nat, Turing.MultiTapeTM.runFrom, Turing.MultiTapeTM.runFrom_succ_eq_step, Turing.MultiTapeTM.runFrom_zero, Turing.MultiTapeTM.step, Turing.moveInputPos_zero}
\difficulty{4}
For every input, every decoded word $u$, every position and every $j \le |u|$: from the rewind state with work head at cell $j - 1$, exactly $j+1$ steps of the pair comparator reach the comparison state with work head at cell zero and all data unchanged.
\end{lemma}

\begin{lemma}[Prefix-equality invariant of the comparison scan]
\lean{Complexity.a3n_eq_scan}
\label{Complexity.a3n_eq_scan}
\leanok
\uses{Complexity.a3nEqCfg, Complexity.a3nEqTM, Complexity.e3c_take_succ_eq, Turing.Action.apply, Turing.Cfg.inputSymbol, Turing.Cfg.workTapeSymbols, Turing.FinTM.bufferTape_nat, Turing.MultiTapeTM.runFrom, Turing.MultiTapeTM.runFrom_succ_eq_step', Turing.MultiTapeTM.step, Turing.inputSymbolInner, Turing.moveInputPos, Turing.moveInputPos_pos_of_ne_right}
\difficulty{5}
Suppose the comparator input has the form $y = p\,v$ and let $u$ be the decoded word. For every $j \le |v|$, after $j$ steps from the comparison state at position $|p|+1$, the comparator is at position $|p|+j+1$ with its flag recording whether the length-$j$ prefixes of $u$ and $v$ agree.
\end{lemma}

\begin{lemma}[Characterization of whole-word equality]
\lean{Complexity.a3n_eq_end}
\label{Complexity.a3n_eq_end}
\leanok
\difficulty{3}
For all binary strings $u$ and $v$: the length-$|v|$ prefix of $u$ equals $v$ and $u$ has no bit at index $|v|$, if and only if $u = v$.
\end{lemma}

\begin{lemma}[Final verdict of the comparison scan]
\lean{Complexity.a3n_eq_compare}
\label{Complexity.a3n_eq_compare}
\leanok
\uses{Complexity.a3nEqCfg, Complexity.a3nEqTM, Complexity.a3n_eq_end, Complexity.a3n_eq_scan, Turing.Action.apply, Turing.Cfg.inputSymbol, Turing.Cfg.workTapeSymbols, Turing.FinTM.bufferTape_nat, Turing.MultiTapeTM.runFrom, Turing.MultiTapeTM.runFrom_succ_eq_step', Turing.MultiTapeTM.step}
\difficulty{4}
Suppose the comparator input has the form $y = p\,v$ and let $u$ be the decoded word. After $|v| + 1$ steps from the comparison state at position $|p|+1$, the pair comparator has halted and its output is the single bit recording whether $u = v$.
\end{lemma}

\begin{lemma}[Linear-time decision of pair equality]
\lean{Complexity.a3n_eq_computes}
\label{Complexity.a3n_eq_computes}
\leanok
\uses{Complexity.a3nEqCfg, Complexity.a3nEqTM, Complexity.a3n_eq_compare, Complexity.a3n_eq_parse, Complexity.a3n_eq_rewind, Turing.Cfg.init, Turing.FinTM.ComputesInTime, Turing.FinTM.bufferTape, Turing.FinTM.computesInTime_iff, Turing.MultiTapeTM.initCfg, Turing.MultiTapeTM.runFrom_add, Turing.pairEncode}
\difficulty{4}
For all binary strings $u$ and $v$, the pair comparator, on the self-delimiting encoding of the pair $(u,v)$, halts within $3|u| + |v| + 4$ steps with output the single bit recording whether $u = v$.
\end{lemma}

\begin{lemma}[Polynomial-time equality of computed words]
\lean{Complexity.a3n_poly_eq}
\label{Complexity.a3n_poly_eq}
\leanok
\uses{Complexity.PolyTimeComputable, Complexity.a3nEqTM, Complexity.a3nVerifier, Complexity.a3n_comp_at, Complexity.a3n_eq_computes, Complexity.a3n_poly_pair, Turing.FinTM.ComputesInTime.mono, Turing.FinTM.bufferedCompTM, Turing.FinTM.computesInTime_iff, Turing.MultiTapeTM.output_length_le, Turing.pairEncode}
\difficulty{4}
Let $f$ and $g$ be polynomial-time computable functions on binary strings. Then the single-bit function recording whether $f(x) = g(x)$ is polynomial-time computable.
\end{lemma}

\begin{lemma}[Polynomial-time equality with a fixed word]
\lean{Complexity.a3n_poly_fixed_eq}
\label{Complexity.a3n_poly_fixed_eq}
\leanok
\uses{Complexity.PolyTimeComputable, Complexity.a3n_poly_linear, Turing.FinTM.computesFunInTime_ifEq}
\difficulty{3}
For every fixed binary word $w$, the single-bit function recording whether the input equals $w$ is polynomial-time computable.
\end{lemma}

\begin{lemma}[A fixed-width increment never returns the empty word]
\lean{Complexity.a3n_inc_nonempty}
\label{Complexity.a3n_inc_nonempty}
\leanok
\uses{Turing.incFixed}
\difficulty{2}
For every binary word $w$, a successful fixed-width increment of $w$ is never the empty word.
\end{lemma}

\begin{lemma}[Increment overflow characterizes the all-ones shape]
\lean{Complexity.a3n_inc_none}
\label{Complexity.a3n_inc_none}
\leanok
\uses{Turing.incFixed}
\difficulty{3}
For every binary word $w$, the fixed-width increment of $w$ overflows if and only if $w$ is the all-ones word of its own length, the empty word included.
\end{lemma}

\begin{lemma}[Polynomial-time test for the all-ones shape]
\lean{Complexity.a3n_poly_unary}
\label{Complexity.a3n_poly_unary}
\leanok
\uses{Complexity.PolyTimeComputable, Complexity.PolyTimeComputable.comp, Complexity.a3n_inc_none, Complexity.a3n_inc_nonempty, Complexity.a3n_poly_fixed_eq, Complexity.a3n_poly_linear, Turing.FinTM.computesFunInTime_incFixed, Turing.incFixed}
\difficulty{4}
The single-bit function recording whether the input is the all-ones word of its own length is polynomial-time computable.
\end{lemma}

\begin{lemma}[Inversion of the pair decoder]
\lean{Complexity.a3n_pair_inverse}
\label{Complexity.a3n_pair_inverse}
\leanok
\uses{Turing.pairDecode, Turing.pairEncode}
\difficulty{5}
For every binary string $z$ and all strings $a, b$: if the self-delimiting pair decoder returns the pair $(a, b)$ on $z$, then $z$ is exactly the pair encoding of $(a, b)$.
\end{lemma}

\begin{definition}[Verifier for the padded language]
\lean{Complexity.a3nVerifier}
\label{Complexity.a3nVerifier}
\leanok
\uses{Turing.pairDecode}
Given natural numbers $C, c$ and a language $V$, the paired verifier language consists of the strings that decode as a pair $(x', u)$ whose first component decodes in turn as a pair $(x, p)$, such that $p$ is the all-ones word of length exactly $C\cdot 2^{(|x|+1)^c}$, the witness $u$ has exactly that same length, and the concatenation $x\,u$ belongs to $V$; strings failing either decoding are excluded.
\end{definition}

\begin{lemma}[Shape and length determine the pad]
\lean{Complexity.a3n_pad_eq}
\label{Complexity.a3n_pad_eq}
\leanok
\difficulty{2}
For every binary word $p$ and natural number $n$: $p$ is the all-ones word of its own length and $|p| = n$, if and only if $p$ is the all-ones word of length $n$.
\end{lemma}

\begin{lemma}[The padded verifier is polynomial time]
\lean{Complexity.a3n_verifier_mem_P}
\label{Complexity.a3n_verifier_mem_P}
\leanok
\uses{Complexity.P, Complexity.PolyTimeComputable, Complexity.PolyTimeComputable.comp, Complexity.a3nVerifier, Complexity.a3n_concat, Complexity.a3n_fst, Complexity.a3n_pad_eq, Complexity.a3n_poly_cond, Complexity.a3n_poly_const, Complexity.a3n_poly_eq, Complexity.a3n_poly_linear, Complexity.a3n_poly_pair, Complexity.a3n_poly_unary, Complexity.a3n_snd, Complexity.e3_exp_bits_timed, Complexity.e3c_bits_injective, Complexity.mem_P_iff, Turing.FinTM.computesFunInTime_lengthBits, Turing.FinTM.computesFunInTime_pairConcat, Turing.FinTM.computesFunInTime_pairFst, Turing.FinTM.computesFunInTime_pairSnd, Turing.FinTM.computesFunInTime_pairValid, Turing.MultiTapeTM.indicator, Turing.pairDecode, Turing.pairDecode_pairEncode}
\difficulty{6}
For all natural numbers $C, c$ and every language $V$ in $\mathrm{P}$, the paired verifier language for $C$, $c$ and $V$, which checks both pair decodings, the exact all-ones pad of length $C\cdot 2^{(|x|+1)^c}$, the exact witness length, and membership of the assembled word in $V$, belongs to $\mathrm{P}$.
\end{lemma}

\begin{definition}[The exponentially padded language]
\lean{Complexity.a3nPadLanguage}
\label{Complexity.a3nPadLanguage}
\leanok
\uses{Turing.pairEncode}
Given a language $L$ and natural numbers $C, c$, the padded language consists of the pair encodings of $(x, p)$ where $x \in L$ and $p$ is the all-ones word of length $C \cdot 2^{(|x|+1)^c}$.
\end{definition}

\begin{lemma}[Padded membership reflects original membership]
\lean{Complexity.a3n_pad_member}
\label{Complexity.a3n_pad_member}
\leanok
\uses{Complexity.a3nPadLanguage, Turing.pairEncode, Turing.pairEncode_injective}
\difficulty{3}
Let $L$ be a language and $C, c$ natural numbers. For every binary string $x$, the pair encoding of $x$ with the all-ones pad of length $C\cdot 2^{(|x|+1)^c}$ belongs to the exponentially padded language of $L$ if and only if $x \in L$.
\end{lemma}

\begin{lemma}[The padded language is in NP]
\lean{Complexity.a3n_pad_mem_NP}
\label{Complexity.a3n_pad_mem_NP}
\leanok
\uses{Complexity.NP, Complexity.P, Complexity.a3nPadLanguage, Complexity.a3nVerifier, Complexity.a3n_pair_inverse, Complexity.a3n_verifier_mem_P, Complexity.mem_NP_iff_exists_length_le, Turing.pairDecode, Turing.pairDecode_pairEncode, Turing.pairEncode}
\difficulty{5}
Let $L$ and $V$ be languages with $V \in \mathrm{P}$, and let $C, c$ be natural numbers such that $x \in L$ if and only if some $u$ with $|u| = C \cdot 2^{(|x|+1)^c}$ has $x\,u \in V$. Then the exponentially padded language of $L$ belongs to $\mathrm{NP}$.
\end{lemma}

\begin{lemma}[Length bounds before validity checks]
\lean{Complexity.a3n_prevalidation}
\label{Complexity.a3n_prevalidation}
\leanok
\uses{Complexity.a3n_pair_inverse, Complexity.e3_bits_length_bound, Turing.pairDecode, Turing.pairEncode}
\difficulty{3}
Let $C, c$ be natural numbers and suppose a string $w$ decodes as a pair with first component $x'$, which in turn decodes as a pair with first component $x$. Then $|x| \le |w|$, and the binary expansion of $C \cdot 2^{(|x|+1)^c}$ has length at most $(|w|+1)^c$ plus the length of the binary expansion of $C$.
\end{lemma}

\begin{lemma}[Envelope for a rescaled exponential width]
\lean{Complexity.a3n_scaled_envelope}
\label{Complexity.a3n_scaled_envelope}
\leanok
\uses{Complexity.a2_exponent_bound}
\difficulty{4}
For all natural numbers $C, c, K, r$ and $s$ there exist constants $A$ and $e$ such that $K\,(n + C\cdot 2^{(s(n+1))^c} + 1)^r \le A \cdot 2^{n^e}$ for every natural number $n$.
\end{lemma}

\begin{lemma}[Emitting the padded pair]
\lean{Complexity.a3n_pad_emit}
\label{Complexity.a3n_pad_emit}
\leanok
\uses{Complexity.a2_exp_scheduler, Complexity.a3n_comp_at, Turing.FinTM, Turing.FinTM.ComputesFunInTime, Turing.FinTM.ComputesInTime.mono, Turing.FinTM.bufferedCompTM, Turing.FinTM.computesFunInTime_pairDup, Turing.FinTM.computesFunInTime_pairMapSnd, Turing.pairDecode_pairEncode, Turing.pairEncode}
\difficulty{5}
For all natural numbers $C$ and $c$ there exist a deterministic machine and a constant $K$ computing, on every input $x$, the pair encoding of $x$ with the all-ones pad of length $C \cdot 2^{(|x|+1)^c}$, within $K\,(n + C\cdot 2^{(3(n+1))^c} + 1)^{c+2}$ steps on inputs of length $n$.
\end{lemma}

\begin{lemma}[Unpadding a polynomial decider into EXP]
\lean{Complexity.a3n_unpad_EXP}
\label{Complexity.a3n_unpad_EXP}
\leanok
\uses{Complexity.EXP, Complexity.P, Complexity.a3nPadLanguage, Complexity.a3n_comp_at, Complexity.a3n_pad_emit, Complexity.a3n_pad_member, Complexity.a3n_scaled_envelope, Complexity.mem_P_iff, Turing.FinTM.ComputesInTime.mono, Turing.FinTM.bufferedCompTM, Turing.MultiTapeTM.indicator, Turing.pairEncode}
\difficulty{5}
Let $L$ be a language over the binary alphabet and $C, c$ natural numbers. If the exponentially padded language of $L$ belongs to $\mathrm{P}$, then $L$ belongs to $\mathrm{EXP}$.
\end{lemma}

\begin{theorem}[P = NP collapses EXP to NEXP]
\lean{Complexity.EXP_eq_NEXP_of_P_eq_NP}
\label{Complexity.EXP_eq_NEXP_of_P_eq_NP}
\leanok
\uses{Complexity.EXP, Complexity.EXP_subset_NEXP, Complexity.NEXP, Complexity.NP, Complexity.P, Complexity.a3nPadLanguage, Complexity.a3n_pad_mem_NP, Complexity.a3n_unpad_EXP}
\difficulty{2}
If $\mathrm{P} = \mathrm{NP}$, then $\mathrm{EXP} = \mathrm{NEXP}$.
\end{theorem}

\begin{theorem}[Separations scale down by padding]
\lean{Complexity.P_ne_NP_of_EXP_ne_NEXP}
\label{Complexity.P_ne_NP_of_EXP_ne_NEXP}
\leanok
\uses{Complexity.EXP, Complexity.EXP_eq_NEXP_of_P_eq_NP, Complexity.NEXP, Complexity.NP, Complexity.P}
\difficulty{1}
If $\mathrm{EXP} \ne \mathrm{NEXP}$, then $\mathrm{P} \ne \mathrm{NP}$.
\end{theorem}
